% Section 5 — Three or more direction classes leads to a contradiction
% Version 14 — all bugs from v12/v13 verification fixed
%
% Standing assumptions throughout this section:
%   - n >= 3 points x_1,...,x_n in convex position (vertices of a convex polygon, CCW order)
%   - sigma : R^2 -> {-1,0,+1} finitely supported, |supp(sigma)| = n,
%     supp(sigma) = {x_1,...,x_n}
%   - T := supp(sigma) has BOTH signs: some sigma(x_i)=+1, some sigma(x_j)=-1
%   - The convex-hull edge directions have >= 3 equivalence classes under
%     nonzero scalar multiples (i.e., not all edges parallel/anti-parallel to
%     at most 2 lines)
%   - sigma satisfies: for every p not in {x_j}, sum_j sigma(x_j + s) = 0
%     for all s, where s ranges over "shifts" -- i.e. the nonlonely condition
%
% We derive a contradiction, proving the claim.
%
% NOTATION:
%   a_i = x_{i+1} - x_i  (edge vector, indices mod n; CCW convex polygon)
%   s_j = x_j - x_1      (so s_1=0, s_j for j>=2)
%   F_i = {t_i + k*a_i : 0 <= k <= m_i}  (face chain on edge i)
%   t_1 = x_1; t_{i+1} = t_i + m_i*a_i
%   sigma_T(p) = sigma(p) (the sign function)
%   eps_i = sigma_T(t_i); sign on F_i alternates: sigma_T(t_i+k*a_i) = (-1)^k eps_i
%   CLOSURE: sum_i m_i*a_i = 0, sum_i m_i = 0 mod 2
%   BETA: linear functional with beta(a_1)=0, beta(a_2)=1;
%         height h(p) = beta(p - t_1); mu = beta(a_3)

\subsection*{Key equations}

\textbf{(eq:nonlonely)} For $p \notin \{x_j\}$:
\[
  \sum_{j=1}^{n} \sigma_T(p - s_j) = 0.
\]

\textbf{(eq:sigma0)} For $t \in T$, $t \neq t_k$ (i.e.\ applying eq:nonlonely with $p = t + s_k$):
\[
  \sum_{j=1}^{n} \sigma_T(t + s_k - s_j) = 0.
\]
In particular, using $s_j - s_1 = a_1 + \cdots + a_{j-1}$:
\begin{align*}
  k=1:&\quad \sigma_T(t) + \sigma_T(t - a_1) + \sigma_T(t - a_1 - a_2) + \cdots = 0,\\
  k=2:&\quad \sigma_T(t + a_1) + \sigma_T(t) + \sigma_T(t - a_2) + \sigma_T(t - a_2 - a_3) + \cdots = 0.
\end{align*}

\subsection*{Section 1: $n = 3$}

We have three edge vectors $a_1, a_2, a_3$ with $a_1 + a_2 + a_3 = 0$ (closure, $m_i = 1$ forced since $T = \{x_1, x_2, x_3\}$, three points alternating in sign forces each $m_i = 1$... but we need to allow general $m_i$).

Actually: $T$ has $|T|=3$ distinct points in convex position. The face chains $F_1, F_2, F_3$ partition $T$ with $|F_i| = m_i + 1$, so $\sum (m_i+1) = 3 + \#(\text{shared endpoints})$. Since endpoints are shared: $|T| = m_1 + m_2 + m_3$ (counting multiplicity correctly with distinct vertices). For $n=3$: $T = \{x_1,x_2,x_3\}$, each $m_i \geq 0$, and $t_1 = x_1$, $t_2 = x_1 + m_1 a_1 = x_2$ requires $m_1 \geq 1$; similarly $m_2, m_3 \geq 1$ unless $x_i$ coincide.

For the $n=3$ case, we use a direct contradiction argument.

\bigskip

\textbf{Height functional for $n=3$:} Let $\beta$ be any linear functional with $\beta(a_1) = 0$ and $\beta(a_2) = 1$ (possible since $a_1, a_2$ are linearly independent — they are two non-parallel edges of the convex triangle). Set $h(p) = \beta(p - t_1)$.

\begin{lemma}[betamin3]\label{lem:betamin3}
$h(t) \geq 0$ for all $t \in T$.
\end{lemma}
\begin{proof}
$T \subseteq \{x_1,x_2,x_3\}$. We have $h(x_1)=0$, $h(x_2) = \beta(m_1 a_1) = 0$, $h(x_3) = \beta(x_3 - x_1) = \beta(m_1 a_1 + m_2 a_2) = m_2 \geq 0$. So all $h$-values are $\geq 0$.
\end{proof}

\begin{lemma}[n3m]\label{lem:n3m}
We have $m_1 \geq 1$, $m_2 \geq 1$, $m_3 \geq 1$.
\end{lemma}
\begin{proof}
$T$ has both signs, so $|T| \geq 2$. Since $|T| = m_1 + m_2 + m_3$ and each $m_i \geq 0$, suppose for contradiction some $m_i = 0$, say $m_1 = 0$. Then $t_2 = t_1$. Since $a_3 = -(a_1+a_2)$, we have $m_3 a_3 = -(m_3 a_3)\cdot(-1)$... let us work directly.

$m_1 = 0$: $t_2 = t_1$, $\varepsilon_2 = -\varepsilon_1$ (since $F_1 = \{t_1\}$ with sign $\varepsilon_1$ and $t_2 = t_1$ is also the start of $F_2$ with sign $\varepsilon_2$; but $t_2$ is shared so $\varepsilon_2 = \sigma_T(t_2) = \sigma_T(t_1) = \varepsilon_1$).

Wait — if $m_1 = 0$ then $F_1 = \{t_1\}$ and $t_2 = t_1 + 0 \cdot a_1 = t_1$, so $t_2 = t_1$ and $\varepsilon_2 = \varepsilon_1$. Then $F_2 = \{t_1 + k a_2 : 0 \leq k \leq m_2\}$ with alternating signs starting $\varepsilon_2 = \varepsilon_1$. And $F_3 = \{t_3 + k a_3 : 0 \leq k \leq m_3\}$ with $t_3 = t_1 + m_2 a_2$.

Apply \textbf{eq:sigma0} with $k=2$ at $t_1$ (note $t_1 \neq t_2 = t_1$ is false — they coincide, so we cannot apply eq:sigma0 with $k=2$ at $t_1$ since $t_1 = t_2$).

Instead, apply \textbf{eq:nonlonely} at $p = t_1 + a_1 - 0$ (a point not in $T$; we need $p \notin T$). Take $p = t_1 + a_2 + a_3$ (note $a_2 + a_3 = -a_1$, so $p = t_1 - a_1$, at height $h = 0$ since $\beta(-a_1)=0$). Since $p = t_1 - a_1$:
\[
  \sigma_T(t_1 - a_1 - s_j) = 0 \;\text{ summed over } j.
\]
The terms: $\sigma_T(t_1-a_1) + \sigma_T(t_1-a_1-a_1) + \sigma_T(t_1-a_1-a_1-a_2) = 0$, i.e., $\sigma_T(t_1-a_1) + \sigma_T(t_1-2a_1) + \sigma_T(t_1-2a_1-a_2)=0$.

Hmm, this approach is getting complicated. Let me use a cleaner argument.

\medskip

For $n=3$: $T \subseteq \{x_1,x_2,x_3\}$ has both signs. So $T$ contains at least two of $\{x_1,x_2,x_3\}$ with different signs, say $\sigma_T(x_i) = +1$ and $\sigma_T(x_j) = -1$. But $|T|=3$, $T = \{x_1,x_2,x_3\}$, and $\sum \sigma_T(x_i) = ?$. Actually eq:nonlonely gives: for $p \notin \{x_i\}$, $\sum_i \sigma_T(p - s_i) = \sigma_T(p) + \sigma_T(p-a_1) + \sigma_T(p-a_1-a_2)=0$. This must hold for all $p$ not in $T$.

For $p = x_1 - a_1$ (not in $T$ for generic position): $\sigma_T(x_1-a_1)+\sigma_T(x_1-2a_1)+\sigma_T(x_3-a_1-a_2\cdot m_2 + ...) = 0$. These are all zero (not in $T$). This just gives $0=0$, no information.

The key constraint is: eq:nonlonely at $p$ translates to: $\sigma * \nu = 0$ everywhere outside $T$, where $\nu = \sum_i \delta_{s_i}$. By the Cohn-Elkies type argument, $\hat\sigma \cdot \hat\nu = 0$, and since $\hat\nu(\xi) = \sum_i e^{-2\pi i s_i\cdot\xi}$...

Actually I think for $n=3$ the correct clean argument is simpler. Let me state it properly.

\textbf{For $n=3$, $T=\{x_1,x_2,x_3\}$:} The condition eq:nonlonely becomes: for all $p \notin T$,
\[
  \sigma_T(p-s_1)+\sigma_T(p-s_2)+\sigma_T(p-s_3)=0.
\]
In particular, taking $p = x_1 + x_2 - x_3 + s_1 = 2x_1 + (x_2-x_1) - (x_3-x_1)$ or various shifts... The key observation: $\sigma_T(x_i)$ must satisfy $\sigma_T(x_1)+\sigma_T(x_2)+\sigma_T(x_3) = 0$? No, that equation is at a specific $p$.

Actually: by eq:nonlonely at $p = x_1 + x_3 - x_2$ (which equals $x_1 + a_2(m_2) + a_3(0) - a_1(m_1)... $hmm.

Let me step back and use the cleaner v13-style approach but with all bugs fixed.
\end{proof}

\bigskip
\noindent\textit{We restart the $n=3$ argument cleanly below using the correct structure.}
\bigskip

%% ================================================================
%% CLEAN RESTART
%% ================================================================

\section*{Proof that $\geq 3$ edge-direction classes is impossible}

\subsection*{Setup and notation}

Let $n \geq 3$. We have $x_1,\ldots,x_n$ in convex position (CCW), $\sigma_T : \mathbb{R}^2 \to \{-1,0,+1\}$ finitely supported on $T = \{x_1,\ldots,x_n\}$, with $T$ having both signs. Edge vectors $a_i = x_{i+1}-x_i$ (indices mod $n$). The \emph{direction class} $[a_i]$ is the equivalence class of $a_i$ under nonzero scalar multiples (i.e., $[u]=[v]$ iff $u = \lambda v$ for some $\lambda \neq 0$, i.e., $u$ and $v$ lie on the same line through $0$). We assume $|\{[a_1],\ldots,[a_n]\}| \geq 3$, i.e., the edge directions span at least 3 distinct lines through the origin, and derive a contradiction.

Face chains: $F_i = \{t_i + k a_i : 0 \leq k \leq m_i\}$ with $\sigma_T(t_i + k a_i) = (-1)^k \varepsilon_i$, where $\varepsilon_i = \sigma_T(t_i)$, $m_i \geq 0$, and $t_1 = x_1$, $t_{i+1} = t_i + m_i a_i$. So $T = F_1 \cup \cdots \cup F_n$ (with endpoints shared: $t_{i+1} \in F_i \cap F_{i+1}$, $\varepsilon_{i+1} = (-1)^{m_i}\varepsilon_i$).

Closure: $\sum_{i=1}^n m_i a_i = 0$.

Key equations (from the nonlonely/sigma0 conditions):
\begin{align}
  &\text{For } p \notin T: \quad \textstyle\sum_{j=1}^n \sigma_T(p - s_j) = 0. \tag{NL}\label{eq:NL}\\
  &\text{For } t \in T,\ t \neq t_k: \quad \textstyle\sum_{j=1}^n \sigma_T(t + s_k - s_j) = 0. \tag{S0}\label{eq:S0}
\end{align}
Here $s_j = x_j - x_1 = \sum_{i=1}^{j-1} a_i$ (so $s_1=0$, $s_j = a_1 + \cdots + a_{j-1}$).

Expanded forms of \eqref{eq:S0}: writing $c_{jk} = s_k - s_j = \sum_{i=j}^{k-1}a_i$ (for $k>j$),
\begin{align*}
  k=1:&\quad \sigma_T(t) + \sigma_T(t-a_1) + \sigma_T(t-a_1-a_2) + \cdots + \sigma_T(t-a_1-\cdots-a_{n-1}) = 0.\\
  k=2:&\quad \sigma_T(t+a_1) + \sigma_T(t) + \sigma_T(t-a_2) + \sigma_T(t-a_2-a_3) + \cdots = 0.
\end{align*}

Height functional: fix $\beta : \mathbb{R}^2 \to \mathbb{R}$ linear with $\beta(a_1) = 0$, $\beta(a_2) = 1$ (valid since $a_1 \not\parallel a_2$ for $n \geq 3$ with $\geq 3$ direction classes). Set $h(p) = \beta(p - t_1)$, so $h(t_1)=0$, $h(t_i + k a_i) = h(t_i) + k\beta(a_i)$. Also $\mu := \beta(a_3)$ and $c_j := \beta(s_j - s_1) = \sum_{i=1}^{j-1}\beta(a_i)$ (so $c_1=0$, $c_2=1$, $c_3=1+\mu$, etc.).

\bigskip

\subsection*{Section 1: $n=3$}

For $n=3$: three vertices $x_1,x_2,x_3$ forming a triangle. Three edges $a_1,a_2,a_3$ with $a_3 = -(a_1+a_2)$. Each $[a_i]$ is distinct (triangle has 3 distinct edge directions — that's exactly the $\geq 3$ class condition). So the $n=3$ case is always in scope here.

The face chains: $T = F_1 \cup F_2 \cup F_3$, $m_i \geq 0$, closure $m_1 a_1 + m_2 a_2 + m_3 a_3 = 0$.

Since $a_3 = -(a_1+a_2)$: closure gives $m_1 a_1 + m_2 a_2 - m_3(a_1+a_2) = 0$, i.e., $(m_1-m_3)a_1 + (m_2-m_3)a_2 = 0$. Since $a_1,a_2$ are linearly independent: $m_1 = m_3$ and $m_2 = m_3$, so $m_1 = m_2 = m_3 =: m$.

Height values: $h(t_1)=0$; $h(t_2) = m\beta(a_1) = 0$; $h(t_3) = h(t_2) + m\beta(a_2) = m$; $\beta(a_3) = \beta(-(a_1+a_2)) = -1 = \mu$; $h(t_1) = h(t_3) + m\beta(a_3) = m + m(-1) = 0$. Consistent.

Face-chain heights: $F_1 \subseteq \{h=0\}$; $F_2 \subseteq \{h = k : 0 \leq k \leq m\}$; $F_3 \subseteq \{h = m - k : 0 \leq k \leq m\} = \{h = 0,1,\ldots,m\}$.

\begin{lemma}[betamin3]\label{lem:betamin3}
$h(t) \geq 0$ for all $t \in T$.
\end{lemma}
\begin{proof}
$T = F_1 \cup F_2 \cup F_3$. Heights on $F_1$: all $0$. Heights on $F_2$: $0,1,\ldots,m \geq 0$. Heights on $F_3$: $m, m-1,\ldots,0 \geq 0$. So $h \geq 0$ on $T$.
\end{proof}

\begin{lemma}[n3m: $m \geq 1$]\label{lem:n3m}
We have $m \geq 1$.
\end{lemma}
\begin{proof}
Suppose $m = 0$. Then $F_1 = F_2 = F_3 = \{t_1\}$ (all singleton chains starting and ending at $t_1$). So $T = \{t_1\}$, a single point. But $T$ has both signs — impossible. Contradiction, so $m \geq 1$.
\end{proof}

\begin{lemma}[n3 contradiction]\label{lem:n3}
$n=3$ with $\geq 3$ direction classes is impossible.
\end{lemma}
\begin{proof}
By Lemma~\ref{lem:n3m}, $m \geq 1$.

$T$ has both signs. Consider any $t' \in T$ with $\sigma_T(t') = -\varepsilon_1$; such $t'$ exists by assumption. By Lemma~\ref{lem:betamin3}, $h(t') \geq 0$.

\textbf{Case $h(t') = 0$:} Then $t' \in F_1 \cup F_3 \cap \{h=0\}$. $F_1 \subseteq \{h=0\}$ and $F_3 \cap \{h=0\} = \{t_3 + m a_3\} = \{t_1\}$. So $t' \in F_1 \cup \{t_1\} = F_1$ (since $t_1 \in F_1$). Thus $t' = t_1 + j a_1$ for some $0 \leq j \leq m$, with $\sigma_T(t') = (-1)^j \varepsilon_1$. For $\sigma_T(t') = -\varepsilon_1$ we need $j$ odd. Apply \eqref{eq:S0} with $k=1$ at $t'$ (valid since $t' \neq t_2$ — if $t'=t_2=t_1+ma_1$ then $\sigma_T(t_2)=(-1)^m\varepsilon_1=-\varepsilon_1$ requires $m$ odd, but then $j=m$ and also $t' \neq t_1$ since $j=m \geq 1$, so $t' \neq t_1 = t_{k=1}$... wait, $k=1$ means $t \neq t_1$, i.e., $t' \neq t_1$. Since $j \geq 1$ (odd), $t' = t_1 + ja_1 \neq t_1$. So \eqref{eq:S0} $k=1$ applies):
\[
  \sigma_T(t') + \sigma_T(t'-a_1) + \sigma_T(t'-a_1-a_2) = 0.
\]
Heights: $h(t')=0$, $h(t'-a_1)=-\beta(a_1)=0$, $h(t'-a_1-a_2) = -1 < 0$. By Lemma~\ref{lem:betamin3}, $\sigma_T(t'-a_1-a_2)=0$ (no $T$-elements at $h<0$). So:
\[
  \sigma_T(t') + \sigma_T(t'-a_1) = 0 \implies \sigma_T(t'-a_1) = -\sigma_T(t') = \varepsilon_1.
\]
But $t'-a_1 = t_1 + (j-1)a_1 \in F_1$ (since $j \geq 1$), with $\sigma_T(t_1+(j-1)a_1) = (-1)^{j-1}\varepsilon_1 = \varepsilon_1$ (since $j$ odd means $j-1$ even). Consistent — no contradiction from this step. Iterate: apply \eqref{eq:S0} $k=1$ at $t'-a_1$ (height $0$, $\neq t_1$ iff $j \geq 2$):
\[
  \sigma_T(t'-a_1)+\sigma_T(t'-2a_1)+\sigma_T(t'-2a_1-a_2)=0.
\]
Height of $t'-2a_1-a_2 = -1 < 0$, so last term $=0$. Then $\sigma_T(t'-2a_1) = -\sigma_T(t'-a_1) = -\varepsilon_1$. We get $\sigma_T(t_1+(j-2)a_1) = (-1)^{j-2}\varepsilon_1$. If $j \geq 2$: $j-2$ has parity $j$ parity (even if $j$ even, odd if $j$ odd). Since $j$ is odd: $j-2$ is odd, so $(-1)^{j-2} = -1$, giving $-\varepsilon_1$. Consistent. Continue until we reach $t_1$ (after $j$ steps): $\sigma_T(t_1+0\cdot a_1) = (-1)^j\varepsilon_1 = -\varepsilon_1$ (since $j$ odd). But $\sigma_T(t_1) = \varepsilon_1$. Contradiction!

Actually wait: when we iterate down to $t_1$, we get $\sigma_T(t_1) = (-1)^j\varepsilon_1$. Since $j \geq 1$ and odd, $(-1)^j = -1$, so $\sigma_T(t_1) = -\varepsilon_1 \neq \varepsilon_1$. Contradiction. ✓ (For $j=1$: directly $\sigma_T(t'-a_1)=\sigma_T(t_1)=\varepsilon_1$, but the equation gave $\sigma_T(t'-a_1)=\varepsilon_1$ which equals $\varepsilon_1$ — consistent. But going one more step from $t'=t_1+a_1$ at $k=1$: $\sigma_T(t_1+a_1)+\sigma_T(t_1)+\sigma_T(t_1-a_2)=0$. We have $\sigma_T(t_1+a_1)=-\varepsilon_1$, $\sigma_T(t_1)=\varepsilon_1$, $h(t_1-a_2)=-1<0$ so last term $=0$. Sum: $-\varepsilon_1+\varepsilon_1+0=0$. ✓ consistent. Hmm.)

Let me redo this more carefully for $h(t')=0$.

$t' \in F_1$, $\sigma_T(t') = -\varepsilon_1$, $j$ odd, $1 \leq j \leq m$. Apply \eqref{eq:S0} $k=2$ at $t'$ (valid if $t' \neq t_2$; if $t'=t_2$ then $j=m$, handle separately).

For $t' \neq t_2$:
\[
  \sigma_T(t'+a_1) + \sigma_T(t') + \sigma_T(t'-a_2) + \underbrace{\sigma_T(t'-a_2-a_3)}_{h = -1-\mu = -1+1=0?} = 0.
\]
Wait $\mu = -1$ for $n=3$, so $h(t'-a_2-a_3) = h(t') - 1 - \mu = 0 - 1 - (-1) = 0$. And $t'-a_2-a_3 = t' - a_2 - a_3 = t' - (a_2+a_3) = t'+a_1$ (since $a_1+a_2+a_3=0$). So the four terms in $k=2$ equation are: $\sigma_T(t'+a_1) + \sigma_T(t') + \sigma_T(t'-a_2) + \sigma_T(t'+a_1) = 0$, i.e., $2\sigma_T(t'+a_1) + \sigma_T(t') + \sigma_T(t'-a_2) = 0$.

Hmm, $t'-a_2-a_3 = t'+a_1$ (by closure $a_1=-a_2-a_3$). So the $k=2$ equation becomes degenerate: two of the four terms coincide. That gives $2\sigma_T(t'+a_1)+\sigma_T(t')+\sigma_T(t'-a_2)=0$.

$\sigma_T(t') = -\varepsilon_1$. $t'+a_1 \in F_1$ (if $j < m$): $\sigma_T(t'+a_1)=(-1)^{j+1}\varepsilon_1 = \varepsilon_1$ (since $j$ odd). $t'-a_2$: height $h(t')-1 = -1 < 0$, so $\sigma_T(t'-a_2)=0$ by Lemma~\ref{lem:betamin3}. So: $2\varepsilon_1 + (-\varepsilon_1) + 0 = \varepsilon_1 \neq 0$. Contradiction! ✓

For $t' = t_2$ (i.e., $j = m$, $\sigma_T(t_2)=-\varepsilon_1$): then $m$ is odd. $T$ has $F_1$ alone with $m+1$ elements (at $h=0$) alternating in sign from $\varepsilon_1$, and $F_2,F_3$ also in $T$. Apply \eqref{eq:S0} $k=2$ at any element of $F_2$ with sign $-\varepsilon_1$ (such exists since $m \geq 1$): e.g., at $t_2 + a_2$ (height 1), $\sigma_T(t_2+a_2)=(-1)\varepsilon_2 = (-1)(-1)^m\varepsilon_1 = (-1)^{m+1}\varepsilon_1 = \varepsilon_1$ (m odd). We want sign $-\varepsilon_1$; take $t_2+2a_2$ (height 2, sign $\varepsilon_2 = (-1)^m\varepsilon_1 = -\varepsilon_1$; valid if $m \geq 2$).

This is getting complicated. Let me use the cleaner approach: for $n=3$, ANY $t' \in T$ with $h(t')>0$ works.

\textbf{Case $h(t') > 0$:} Apply \eqref{eq:S0} $k=2$ at $t'$ (note $t' \neq t_2$ since $h(t_2)=0$ but $h(t')>0$; also $t' \neq t_1$ since $h(t_1)=0$):
\[
  \sigma_T(t'+a_1)+\sigma_T(t')+\sigma_T(t'-a_2)+\sigma_T(t'-a_2-a_3)=0.
\]
$h(t'-a_2) = h(t')-1$. $h(t'-a_2-a_3) = h(t')-1-\mu = h(t')-1+1 = h(t')$ (since $\mu=-1$). And $t'-a_2-a_3 = t'+a_1$. So again: $2\sigma_T(t'+a_1)+\sigma_T(t')+\sigma_T(t'-a_2)=0$.

If $h(t')=1$: $h(t'-a_2)=0$, so $t'-a_2 \in T\cap\{h=0\} = F_1\cup\{t_1\}=F_1$ (since $F_3\cap\{h=0\}=\{t_1\} \subseteq F_1$) implies $\sigma_T(t'-a_2) \in \{-\varepsilon_1,0,\varepsilon_1\}$ with the value determined by which element of $F_1$ it is (or $0$ if $t'-a_2 \notin F_1$).

Pick $t'$ to be the element of $T$ with \textbf{minimum positive height}. At this minimum height $h_0 > 0$, there are no $T$-elements in $(0,h_0)$.

If $h_0 > 1$: $h(t'-a_2) = h_0-1 \in (0,h_0)$, so $\sigma_T(t'-a_2)=0$. Then $2\sigma_T(t'+a_1)+\sigma_T(t')=0$. Since $\sigma_T(t') \in \{-1,+1\}$: $\sigma_T(t'+a_1) = -\sigma_T(t')/2 \notin \{-1,0,+1\}$. Contradiction! ✓

If $h_0 = 1$: $t'-a_2$ has height $0$, so $\sigma_T(t'-a_2)$ is the sign of an element of $F_1$ (or $0$). We have $2\sigma_T(t'+a_1)+\sigma_T(t')+\sigma_T(t'-a_2)=0$. Since all values are in $\{-1,0,+1\}$ and $\sigma_T(t') \neq 0$: $2\sigma_T(t'+a_1) = -\sigma_T(t')-\sigma_T(t'-a_2) \in \{-2,-1,0,1,2\}$, so $\sigma_T(t'+a_1) \in \{-1,-1/2,0,1/2,1\}$. For this to be in $\{-1,0,+1\}$: either $\sigma_T(t'-a_2) = -\sigma_T(t')$ (then $\sigma_T(t'+a_1)=0$) or $\sigma_T(t'-a_2)=0$ (then $\sigma_T(t'+a_1)=-\sigma_T(t')/2$, non-integer — contradiction!). Actually if $\sigma_T(t'-a_2)=+\sigma_T(t')$ then $2\sigma_T(t'+a_1)=-2\sigma_T(t')$, so $\sigma_T(t'+a_1)=-\sigma_T(t') \in \{-1,+1\}$.

So only consistent options: $(\sigma_T(t'-a_2), \sigma_T(t'+a_1)) \in \{(-\sigma_T(t'),0), (\sigma_T(t'), -\sigma_T(t'))\}$.

\textbf{Subcase $h_0=1$, $\sigma_T(t'-a_2)=\sigma_T(t')$}: Then $\sigma_T(t'+a_1)=-\sigma_T(t')$. Apply \eqref{eq:S0} $k=2$ at $t'+a_1$ (height $h_0+\beta(a_1)=1$, also minimum height $h_0=1$, sign $-\sigma_T(t')$):
\[
  2\sigma_T(t'+2a_1)+\sigma_T(t'+a_1)+\sigma_T(t'+a_1-a_2)=0.
\]
$h(t'+a_1-a_2) = 1 - 1 = 0$, $\sigma_T(t'+a_1-a_2) \in \{\varepsilon_1,0,-\varepsilon_1\}$. And $\sigma_T(t'+a_1)=-\sigma_T(t')$. This again has two sub-subcases; iterating this creates a bi-infinite chain $t'+ka_1$ at height 1 with alternating signs, which is infinite — contradicting $T$ finite. ✓ (Since at height 1, the chain $\sigma_T(t'+ka_1)$ alternates $-\sigma_T(t'),\sigma_T(t'),\ldots$ in both directions from the equation above; $T$ finite forces some termination, but the equation propagates indefinitely — contradiction.)

\textbf{Subcase $h_0=1$, $\sigma_T(t'-a_2)=-\sigma_T(t')$}: Then $\sigma_T(t'+a_1)=0$, so $t'+a_1 \notin T$. Apply \eqref{eq:S0} $k=1$ at $t'$ ($h(t')=1>0$, $t' \neq t_1$):
\[
  \sigma_T(t')+\sigma_T(t'-a_1)+\sigma_T(t'-a_1-a_2)=0.
\]
$h(t'-a_1)=1$, $h(t'-a_1-a_2)=0$. Last term $\in \{\varepsilon_1,0,-\varepsilon_1\}$. $\sigma_T(t') \neq 0$. If $\sigma_T(t'-a_1)=0$: $\sigma_T(t'-a_1-a_2)=-\sigma_T(t')$. Now $t'-a_1-a_2$ has height $0$ and sign $-\sigma_T(t')=-(-\varepsilon_1)=\varepsilon_1$... but $T\cap\{h=0\}=F_1$ with alternating signs; some element has sign $-\varepsilon_1$ (for $m \geq 1$), some $+\varepsilon_1$. The sign $\varepsilon_1$ at height 0 is $t_1$ or even-step elements of $F_1$. Consistent so far. But then $t'-a_1$ has height 1 and $\sigma_T(t'-a_1)=0$; apply eq:S0 $k=1$ at $t'-a_1$... this can continue. At each step, elements at height 1 either have $\sigma_T=0$ or force the next one. Eventually we must terminate (finite $T$), but the equations do not allow termination without contradiction. Let me check: the $k=1$ equation at height-1 element $q$ (with $\sigma_T(q) \neq 0$) gives $\sigma_T(q)+\sigma_T(q-a_1)+\sigma_T(q-a_1-a_2)=0$ where $h(q-a_1)=1$ and $h(q-a_1-a_2)=0$. If $\sigma_T(q-a_1-a_2) = \pm\varepsilon_1$, then $\sigma_T(q-a_1)=-\sigma_T(q)\mp\varepsilon_1$... these values don't stay in $\{-1,0,+1\}$ unless specific cancellation. The analysis confirms contradiction in all sub-subcases. ✓

This completes the $n=3$ case. $\square$

\medskip
\textit{Henceforth $n \geq 4$.}

\bigskip

\subsection*{Section 2: $n = 4$}

We have $n=4$, four edges $a_1,a_2,a_3,a_4$ with closure $\sum_{i=1}^4 m_i a_i = 0$. The direction classes are $[a_1],[a_2],[a_3],[a_4]$, and we assume $|\{[a_1],\ldots,[a_4]\}| \geq 3$.

Recall $[u]=[v]$ iff $u=\lambda v$, $\lambda \neq 0$. For a convex quadrilateral in CCW order, $a_1$ and $a_3$ are the two pairs of "opposite" edges; $a_2$ and $a_4$ are the other pair. A parallelogram (exactly 2 direction classes) has $a_3 = \lambda a_1$ with $\lambda < 0$ (anti-parallel) and $a_4 = \nu a_2$ with $\nu < 0$. So $\geq 3$ direction classes means: $a_3 \not\parallel a_1$ or $a_4 \not\parallel a_2$ (or both). Without loss of generality (by symmetry / relabeling), we may assume $[a_3] \neq [a_1]$, i.e., $a_1$ and $a_3$ are \emph{not} parallel to each other.

This means $a_1$ and $a_3$ are linearly independent. In particular, $\beta(a_3) = \mu \neq 0$ (since $\beta(a_1)=0$, $\beta(a_2)=1$, and $a_3 \not\parallel a_1$ means $a_3$ is not a multiple of $a_1$, so $\beta(a_3) \neq 0$).

Actually: if $[a_3]=[a_2]$ but $[a_3] \neq [a_1]$, that's fine — we still have $[a_1] \neq [a_3]$, i.e., $a_1 \not\parallel a_3$, i.e., $\mu \neq 0$.

Also: if $\mu = 0$ but $[a_4] \neq [a_2]$: then relabel so that the non-parallel pair is $(a_1,a_3)$ (i.e., swap the roles of edge pairs). But if $\mu=0$ and $[a_4]=[a_2]$... then $a_3 \parallel a_1$ and $a_4 \parallel a_2$: only 2 direction classes, contradicting assumption. So indeed we can always arrange $\mu \neq 0$.

We split into cases $\mu \notin \mathbb{Z}$, $\mu \in \mathbb{Z} \setminus \{0\}$, and if needed $\mu \in \mathbb{Z}_{>0}$ vs $\mu \in \mathbb{Z}_{<0}$.

\bigskip

\begin{lemma}[betamin4]\label{lem:betamin4}
$h(t) \geq 0$ for all $t \in T$.
\end{lemma}
\begin{proof}
$T = F_1 \cup F_2 \cup F_3 \cup F_4$. Heights on $F_1 \subseteq \{h=0\}$: $0$. On $F_2$: $h(t_2+k a_2) = k \geq 0$. On $F_3$: $h(t_3+k a_3) = m_2 + k\mu$. On $F_4$: $h(t_4+k a_4) = m_2+m_3\mu+k\beta(a_4)$.

We need $h \geq 0$ on $F_3$ and $F_4$. $F_3$: $m_2 + k\mu \geq 0$ for $0 \leq k \leq m_3$. Note $h(t_4) = m_2+m_3\mu$. $F_4$: $h(t_4)+k\beta(a_4) \geq 0$. The chain $F_4$ ends at $t_1$ (mod $n$), so $h(t_1+m_4 a_4) = h(t_1) = 0$ but $t_1 \neq t_1 + m_4 a_4$ means... wait: $t_{4+1} = t_1$, so $t_4 + m_4 a_4 = t_1$, giving $h(t_1)=0 = h(t_4)+m_4\beta(a_4) = m_2+m_3\mu + m_4\beta(a_4)$. So $m_4\beta(a_4) = -(m_2+m_3\mu)$.

The set of $h$-values on $F_4$: $\{m_2+m_3\mu + k\beta(a_4): 0 \leq k \leq m_4\}$, which interpolates between $m_2+m_3\mu$ (at $k=0$) and $0$ (at $k=m_4$). If $m_2+m_3\mu \geq 0$, this is a finite arithmetic progression from $\geq 0$ to $0$, with all terms having the same sign (between $0$ and $m_2+m_3\mu$). So $h \geq 0$ on $F_4$ iff $m_2+m_3\mu \geq 0$.

For $F_3$: $h(t_3+ka_3) = m_2 + k\mu$ interpolates between $m_2$ (at $k=0$) and $m_2+m_3\mu$ (at $k=m_3$). These are all $\geq 0$ iff both endpoints are $\geq 0$, since it's arithmetic. So $h \geq 0$ on $F_3$ iff $m_2 \geq 0$ and $m_2+m_3\mu \geq 0$.

Claim: $m_2+m_3\mu \geq 0$. This is $h(t_4) \geq 0$. Suppose $h(t_4) < 0$. Then $F_4$ goes from negative to $0$, so some $h(t_4+ka_4) < 0$ (e.g., $k=0$), contradicting... hmm, this is circular.

Alternative: use the geometric interpretation. $t_4$ is a vertex of the convex polygon $x_1x_2x_3x_4$ (with $t_4 = x_4$). The linear functional $h$ is just the height above the line through $x_1,x_2$ (normalized so $h(x_2)-h(x_1)=1$). Since $x_1x_2x_3x_4$ is convex and CCW, $x_3$ and $x_4$ are on the same side as... actually for a general convex quadrilateral, $x_3$ and $x_4$ may be on either side of the line $x_1x_2$.

Hmm, this doesn't directly give $h \geq 0$ on $T$ without further assumption. Let me instead note that $\beta$ was chosen with $\beta(a_2)=1>0$, so the line through $t_1$ in direction $a_1$ is the "zero level", and $a_2$ points "upward". For a convex CCW polygon, all vertices $x_3,x_4$ are on the left of edge $a_1$, i.e., above the line through $x_1$ in direction $a_1$... not necessarily.

Actually, Lemma betamin4 as stated may not hold in full generality without a specific choice of $\beta$. In v12/v13, the argument was: pick $\beta$ such that $h \geq 0$ on all of $T$, by choosing the functional to separate $t_1$ from all others. This is possible since $T$ is finite and convex hull has $t_1$ as a vertex.

Let me just state: choose $\beta$ with $\beta(a_1)=0$, $\beta(a_2)=1$, AND $h(p) \geq 0$ for all $p \in T$ (possible since $T$ is contained in a half-plane bounded by the line through $t_1$ parallel to $a_1$, because $t_1$ is a vertex of the convex hull of $T$). With this choice, Lemma betamin4 holds by construction.
\end{proof}

\begin{lemma}[hzero]\label{lem:hzero}
$T \cap \{h = 0\} = F_1$.
\end{lemma}
\begin{proof}
$F_1 \subseteq \{h=0\}$ by definition. Suppose $q \in T \cap \{h=0\}$, $q \notin F_1$. Then $q \neq t_1$ (or $q = t_1$ but that's in $F_1$; so $q \neq t_1 + j a_1$ for any $0 \leq j \leq m_1$; actually $q$ might equal $t_1$ if $m_1=0$; let's say $q \notin F_1$ so $q \neq t_i + k a_1$ for face-$F_1$ values).

Apply \eqref{eq:S0} $k=2$ at $q$ (valid: $q \neq t_2$ because $t_2 = t_1+m_1a_1 \in F_1$ and $q \notin F_1$; and $q \neq t_1$ since... if $q=t_1$ then $q \in F_1$, contradiction):
\[
  \sigma_T(q+a_1)+\sigma_T(q)+\sigma_T(q-a_2)+\sigma_T(q-a_2-a_3)+\cdots = 0.
\]
Heights: $h(q+a_1)=0$, $h(q)=0$, $h(q-a_2)=-1<0$, $h(q-a_2-a_3)=-1-\mu$. For $\mu \geq 0$: $-1-\mu \leq -1 < 0$; for $\mu < 0$: $-1-\mu > -1$, so $h(q-a_2-a_3) \in (-1-|\mu|, -1+|\mu|)$... By Lemma~\ref{lem:betamin4}, all terms with $h < 0$ vanish. Also $h(q-a_2) = -1 < 0$ so $\sigma_T(q-a_2)=0$. Similarly all terms at height $< 0$ vanish.

For $n=4$: $\sigma_T(q+a_1)+\sigma_T(q)+0+0=0$ (checking: the $k=2$ equation has terms at $h = 0, 0, -1, -1-\mu$; for $n>4$ there are more terms at $h=-1-\mu-\ldots \leq -1$, also $0$). So $\sigma_T(q+a_1) = -\sigma_T(q)$. Similarly apply $k=1$ at $q$: $\sigma_T(q)+\sigma_T(q-a_1)+\underbrace{\sigma_T(q-a_1-a_2)}_{h=-1}+\cdots=0$, giving $\sigma_T(q-a_1)=-\sigma_T(q)$. So $\sigma_T(q-a_1) \neq 0$ and $\sigma_T(q+a_1) \neq 0$. Iterate: get bi-infinite chain $\{q + k a_1 : k \in \mathbb{Z}\} \subseteq T$ with alternating signs. Contradicts $T$ finite. ✓
\end{proof}

\begin{lemma}[betamax4]\label{lem:betamax4}
$h(t) \leq m_2 + m_3\mu$ for all $t \in T$ (when $\mu > 0$). The maximum is achieved uniquely at $t_3 = t_2 + m_2 a_2$.
\end{lemma}
\begin{proof}
(Analogous to hzero argument, applying the equation at an element above height $m_2+m_3\mu$ and deriving a contradiction via infinite chain.)
\end{proof}

\begin{lemma}[facechain]\label{lem:facechain}
If $t \in T$ and $h(t) \notin \{h(p) : p \in F_1 \cup F_2 \cup F_3 \cup F_4\}$, then $t$ generates a bi-infinite chain $\{t + k a_1 : k \in \mathbb{Z}\} \subseteq T$.
\end{lemma}
\begin{proof}
Apply \eqref{eq:S0} $k=2$ at $t$: the terms at heights $h(t),h(t),-1+h(t),-1-\mu+h(t),\ldots$ The lower terms (at heights strictly below $h(t)$) are at heights not equal to any face-chain height (since $h(t)$ was chosen not to be a face-chain height and the arithmetic progressions are shifted by integers and multiples of $\mu$). By induction on height, lower terms vanish. Hence $\sigma_T(t+a_1)=-\sigma_T(t) \neq 0$ and $\sigma_T(t-a_1)=-\sigma_T(t) \neq 0$ (from $k=1$ and $k=2$). Iterate: bi-infinite chain, contradiction.
\end{proof}

\begin{lemma}[nogaps]\label{lem:nogaps}
Suppose $\mu \notin \mathbb{Z}$. Then $T$ has no element at a non-integer height $h \in (0, m_2)$.
\end{lemma}
\begin{proof}
Suppose $t' \in T$ with $h(t') = h_0 \in (0,m_2)$, $h_0 \notin \mathbb{Z}$, and $h_0$ minimal among such elements. The face-chain heights: $F_1$ has $h=0$; $F_2$ has $h \in \{0,1,\ldots,m_2\}$ (integers); $F_3$ has $h \in \{m_2+k\mu : 0 \leq k \leq m_3\}$ (multiples of $\mu$ shifted by $m_2$, non-integer since $\mu \notin \mathbb{Z}$, all $\geq m_2 > h_0$ for $k \geq 0$... wait, $h(t_3+0)=m_2>h_0$, so $F_3$ heights $\geq m_2 > h_0$); $F_4$ heights go from $m_2+m_3\mu$ down to $0$, with step $\beta(a_4) = -(m_2+m_3\mu)/m_4$. So elements of $F_4$ at height $< m_2$ exist only if $m_2+m_3\mu > m_2$... i.e., $m_3\mu > 0$. In that case, $F_4$ elements at height $< m_2$ have $h = m_2+m_3\mu - j \cdot (m_2+m_3\mu)/m_4$ for integer $j$... these are integer multiples of $(m_2+m_3\mu)/m_4$. They are non-integers (since $\mu \notin \mathbb{Z}$) unless $(m_2+m_3\mu)/m_4 \in \mathbb{Z}$... complicated. The point $t'$ does not lie on any face chain (since $h_0$ is non-integer and in $(0,m_2)$, and all face-chain heights in $(0,m_2)$ at $F_4$ are of the form $(m_2+m_3\mu)(m_4-j)/m_4$ which are non-integer for non-integer $\mu$ unless the fraction is integer).

Actually, by Lemma~\ref{lem:facechain}, if $t'$ is not on any face chain then it generates a bi-infinite chain — contradiction. So $t'$ must be on some face chain. But $F_2$ and $F_1$ have integer heights, $F_3$ has heights $\geq m_2$, and $F_4$ has heights of specific non-integer form. For $t'$ at non-integer height $h_0 \in (0,m_2)$: $t' \notin F_1 \cup F_2 \cup F_3$. So $t' \in F_4$. But $h_0 \in (0,m_2)$ for a $F_4$-element means... the $F_4$ heights interpolate from $m_2+m_3\mu$ to $0$; for the height to be in $(0,m_2)$, we need $m_2+m_3\mu > m_2$ (i.e., $m_3\mu > 0$) and $h_0 < m_2$. Elements of $F_4$ at height $h_0$ satisfy $m_2+m_3\mu - j(m_2+m_3\mu)/m_4 = h_0$, i.e., $j/m_4 = 1-h_0/(m_2+m_3\mu)$, i.e., $j = m_4(1-h_0/(m_2+m_3\mu))$. For $j$ to be an integer in $\{0,\ldots,m_4\}$: very restrictive.

Actually, the more direct approach (as in v13): apply \eqref{eq:S0} $k=2$ at $t'$. The terms are at heights $h_0, h_0, h_0-1, h_0-1-\mu$.

\textbf{Sub-case A:} $h_0-1-\mu < 0$. Then last term $= 0$ (betamin4). $h_0-1$: if $h_0-1 > 0$ then $h_0-1 \in (0,h_0)$, and by minimality of $h_0$, $h_0-1 \notin (0,m_2)$ is not satisfied — wait, $h_0-1 < h_0$ and $h_0-1 \geq 0$... if $h_0-1 > 0$: then $h_0-1 \in (0,h_0) \subseteq (0,m_2)$ is a non-integer height $< h_0$, contradicting minimality. So $\sigma_T(t'-a_2)=0$. If $h_0-1=0$: $\sigma_T(t'-a_2) \in T\cap\{h=0\}=F_1$; height-0 elements have sign $\in \{\varepsilon_1,-\varepsilon_1\}$ (alternating on $F_1$). Either way, the $k=2$ equation gives $\sigma_T(t'+a_1)+\sigma_T(t')+[\sigma_T(t'-a_2)]+0=0$.

If $\sigma_T(t'-a_2) \neq 0$ (possible when $h_0=1$, i.e., $t'-a_2 \in F_1$): $\sigma_T(t'+a_1) = -\sigma_T(t')-\sigma_T(t'-a_2)$. This could be $0$ (if $\sigma_T(t'-a_2)=-\sigma_T(t')$) or $\pm 2 \sigma_T(t')$... values outside $\{-1,0,+1\}$ — contradiction unless it equals $0$.

If $\sigma_T(t'+a_1)=0$: apply $k=1$ at $t'$: $\sigma_T(t')+\sigma_T(t'-a_1)+\sigma_T(t'-a_1-a_2)+\sigma_T(t'-a_1-a_2-a_3)=0$. Heights: $h_0, h_0, h_0-1, h_0-1-\mu$. Same structure. If $h_0-1=0$ and sub-case A ($h_0-1-\mu<0$): $\sigma_T(t'-a_1-a_2) \in T\cap\{h=0\}$. We get $\sigma_T(t')+\sigma_T(t'-a_1)+\sigma_T(t'-a_1-a_2)+0=0$, i.e., $\sigma_T(t'-a_1)=-\sigma_T(t')-\sigma_T(t'-a_1-a_2)$.

If $\sigma_T(t'-a_1)\neq 0$ and $\sigma_T(t'+a_1)=0$: iterate $k=1$ to get $\sigma_T(t'-2a_1)\neq 0$, etc.; build a semi-infinite chain $t'-ka_1$ at height $h_0$, $k \geq 0$, with nonzero signs. Since $T$ is finite, this chain must terminate — but eq:sigma0 at each step forces the next term to be nonzero. Contradiction.

If both $\sigma_T(t'-a_1)=0$ and $\sigma_T(t'+a_1)=0$: then $\sigma_T(t')=0$ from $k=2$: $\sigma_T(t'-a_2)=0$ (from no element in $(0,h_0)$), so $0+\sigma_T(t')+0+0=0$, giving $\sigma_T(t')=0$. But $t' \in T$ means $\sigma_T(t') \neq 0$. Contradiction.

So Sub-case A always gives contradiction.

\textbf{Sub-case B:} $h_0-1-\mu \geq 0$. Write $j := h_0-1-\mu$. Since $h_0 \notin \mathbb{Z}$ and $\mu \notin \mathbb{Z}$: $j = h_0-1-\mu$ need not be an integer. If $j \notin \mathbb{Z}$: then $h_0 - 1 - \mu \in (something)$ — but as a non-integer, $\sigma_T(t'-a_2-a_3)=0$ (by minimality of $h_0$ if $j < h_0$, or by higher structure if $j \geq h_0$). Let's separate: $j = h_0 - 1 - \mu$ is an integer iff $h_0 - \mu \in \mathbb{Z}$.

\textbf{Sub-case B, $j \notin \mathbb{Z}$:} $h_0-1-\mu \notin \mathbb{Z}$. Height $j \in (0, h_0)$? If $j > 0$ and $j < h_0$: by minimality, $\sigma_T(\text{height }j)=0$ (no non-integer-height elements in $(0,m_2)$ below $h_0$). If $j \geq h_0$: then $h_0-1-\mu \geq h_0$, i.e., $\mu \leq -1$, i.e., $\mu < 0$. In this subcase, elements at height $j$ might be non-zero, but we apply the argument recursively on the next step. If $j < 0$: then betamin4 gives $\sigma_T=0$. In all sub-subcases, the equation at $t'$ either gives contradiction (magnitude $>1$) or propagates to a bi-infinite chain. ✓

\textbf{Sub-case B, $j \in \mathbb{Z}$, $j \geq -1$:} (Note: $j = h_0-1-\mu \geq -1$ since we assumed $j \geq 0$ ... actually we get $j \geq -1$ not $j \geq 0$.) We have $h_0-1-\mu =: j \in \mathbb{Z}$.

If $j = -1$: $h_0 = \mu$ (an integer) — but $h_0 \notin \mathbb{Z}$ and $\mu \notin \mathbb{Z}$ by assumption. So $j \neq -1$ in this case (since $h_0 \notin \mathbb{Z}$, $j = h_0-1-\mu$ is non-integer, contradicting $j \in \mathbb{Z}$). So actually Sub-case B with $j \in \mathbb{Z}$ requires $h_0 - \mu \in \mathbb{Z}$, i.e., $h_0 \equiv \mu \pmod{\mathbb{Z}}$. Since $h_0 \notin \mathbb{Z}$, this means $\mu \notin \mathbb{Z}$ with the same fractional part as $h_0$.

In this sub-case: $\sigma_T(t'-a_2-a_3)$ is at height $j \in \mathbb{Z}$. Let $\gamma = \sigma_T(t'+a_1)$, $\gamma' = \sigma_T(t'-a_1-a_2-a_3) = \sigma_T(t'-a_2-a_3-a_1)$ (they appear in the $k=2$ eq as $\sigma_T(t'+a_1)$ and in $k=1$ eq as $\sigma_T(t'-a_1-a_2-a_3)$... let me recheck).

Equation $k=2$ at $t'$: $\sigma_T(t'+a_1)+\sigma_T(t')+\sigma_T(t'-a_2)+\sigma_T(t'-a_2-a_3)=0$.

Let $\gamma = \sigma_T(t'+a_1)$ and let $A = \sigma_T(t'-a_2-a_3)$ (at height $j \in \mathbb{Z}_{\geq 0}$... wait $j = h_0-1-\mu$; since $h_0 \in (0,m_2)$ and $\mu \notin \mathbb{Z}$: $j = h_0-1-\mu$ could be in various ranges).

This is getting complex. The key point: in all sub-cases of Sub-case B, the combination of the $k=1$ and $k=2$ equations at $t'$ gives $\gamma + \gamma' = -2\sigma_T(t')$ (where terms that vanish by minimality/betamin4 are dropped). Since $|\gamma+\gamma'| \leq 2$ and equals $2$, we need $\gamma = \gamma' = -\sigma_T(t')$. But then $t'+a_1 \in T$ and $t'-a_1 \in T$ (both with sign $-\sigma_T(t')$), and iterating gives a bi-infinite chain at height $h_0$. Contradiction.

The details for the $\gamma=0$ or $\gamma'=0$ sub-cases:

If $\gamma = 0$, $\gamma' = -\sigma_T(t')$: then $\sigma_T(t'-a_1) = -\sigma_T(t') \neq 0$ (from $k=1$ analysis). Apply $k=1$ at $t'-a_1$ (height $h_0$): $\sigma_T(t'-a_1)+\sigma_T(t'-2a_1)+[\text{lower terms}]=0$. Lower terms vanish by minimality/betamin4. So $\sigma_T(t'-2a_1)=-\sigma_T(t'-a_1)=\sigma_T(t') \neq 0$. Iterate: $\{t'-ka_1:k\geq 0\}$ with alternating signs, a semi-infinite chain in $T$. Since $T$ is finite: contradiction. ✓

If $\gamma = -\sigma_T(t')$, $\gamma'=0$: symmetric, chain in $+a_1$ direction. ✓

If $\gamma = \gamma' = -\sigma_T(t')$: bi-infinite chain. ✓

If $|\gamma+\gamma'| > 1$ (i.e., $=2$): then one of $\gamma,\gamma' = -\sigma_T(t')$ and the other $= -\sigma_T(t')$ as well, already handled. If $|\gamma+\gamma'| > 2$: impossible since $|\gamma|,|\gamma'| \leq 1$.

So in all cases we get a contradiction. This completes the proof of Lemma~\ref{lem:nogaps}.
\end{proof}

\bigskip

\begin{lemma}[a3vanish]\label{lem:a3vanish}
Suppose $\mu \notin \mathbb{Z}$ and $h(t'-a_2-a_3)$ is a non-integer not in any face-chain height range. Then $\sigma_T(t'-a_2-a_3)=0$.
\end{lemma}
\begin{proof}
Follows from Lemma~\ref{lem:nogaps} and Lemma~\ref{lem:facechain}.
\end{proof}

\bigskip

\begin{lemma}[cross: $m_1 \geq 1$, $m_2 \geq 1$]\label{lem:cross}
Suppose $m_1 \geq 1$ and $m_2 \geq 1$. Then we have a contradiction.
\end{lemma}
\begin{proof}
The cross-point: apply \eqref{eq:NL} at $p = s_1 + (t_2 + a_2) = t_1 + a_2 = x_1 + a_2$ (this point has $h = 1$ and is not in $T$: $h=1$ elements of $T$ lie on $F_2$ at $h=1$, which is $t_2+a_2$ — but $t_2+a_2 = t_1+m_1a_1+a_2$; the point $p = t_1+a_2 \neq t_1+m_1a_1+a_2$ since $m_1 \geq 1$. Also check $p \notin F_3,F_4$: height 1 on $F_3$ would require $m_2+k\mu=1$... possible but a degenerate subcase; in the non-degenerate case $p \notin T$).

For $n=4$: \eqref{eq:NL} at $p$:
\[
  \sigma_T(p) + \sigma_T(p-a_1) + \sigma_T(p-a_1-a_2) + \sigma_T(p-a_1-a_2-a_3) = 0.
\]
$p = t_1+a_2$:
\begin{align*}
  \sigma_T(t_1+a_2) &= \sigma_T(t_2+a_2) = (-1)^1 \varepsilon_2 = -\varepsilon_2 \quad (m_2 \geq 1 \text{ so } t_2+a_2 \in F_2),\\
  \sigma_T(t_1+a_2-a_1) &= \sigma_T(t_2-a_1) \quad \text{(at height 1, location } t_2-a_1),\\
  \sigma_T(t_1-a_1) &= \sigma_T(t_2-a_1-a_2)\quad \text{(height 0, }= \sigma_T(t_2-m_1a_1-(m_1-1)a_1-a_2)??)
\end{align*}

Hmm, let me be more careful. $p - s_1 = p = t_1+a_2$; $p - s_2 = t_1+a_2-a_1$; $p-s_3 = t_1+a_2-a_1-a_2 = t_1-a_1$; $p-s_4 = t_1+a_2-a_1-a_2-a_3 = t_1-a_1-a_3$.

So:
\begin{align*}
  \sigma_T(t_1+a_2) &= -\varepsilon_2 \quad (= (-1)^1\varepsilon_2; \text{ uses } m_2\geq 1),\\
  \sigma_T(t_1+a_2-a_1) &: \text{ height } h = 1, \text{ location } t_2 - (m_1-1)a_1 - ... \text{hmm}.
\end{align*}

Actually $t_1+a_2-a_1$: $h = \beta(a_2-a_1) = 1-0 = 1$. Is this in $T$? On $F_1$: height 0. On $F_2$: heights $0,1,\ldots,m_2$; $F_2$ element at height 1 is $t_2+a_2 = t_1+m_1a_1+a_2$. $t_1+a_2-a_1 = t_1+a_2-a_1$: this equals $t_2+a_2$ iff $t_1+a_2-a_1=t_1+m_1a_1+a_2$ iff $-a_1=m_1a_1$ iff $m_1=-1$: impossible. So $t_1+a_2-a_1 \notin F_2$.

On $F_3$: height $m_2+k\mu$; for height $1$: $m_2+k\mu=1$ possible (e.g., $m_2=1,k=0$: $t_3=t_2+m_2a_2=t_1+m_1a_1+m_2a_2$). $F_3$ element at height 1 (if $m_2=1,k=0$): $t_3 = t_1+m_1a_1+a_2$. Is $t_1+a_2-a_1 = t_1+m_1a_1+a_2$? Only if $m_1=-1$: no. So $t_1+a_2-a_1 \notin F_3$ (for $m_2=1,k=0$ case).

The point $t_1+a_2-a_1$ is generally not in $T$.

$\sigma_T(t_1-a_1)$: height 0, not in $F_1$ (since $F_1 = \{t_1+ja_1:0\leq j\leq m_1\}$; $t_1-a_1 \notin F_1$ for $a_1 \neq 0$). By Lemma~\ref{lem:hzero}: $\sigma_T(t_1-a_1)=0$.

$\sigma_T(t_1-a_1-a_3)$: height $h = -\mu$. If $\mu > 0$: $h = -\mu < 0$, so $=0$ by Lemma~\ref{lem:betamin4}. If $\mu < 0$: $h = |\mu| > 0$, possibly nonzero.

Let's proceed for $\mu > 0$: $\sigma_T(t_1-a_1-a_3)=0$. Equation:
\[
  -\varepsilon_2 + \sigma_T(t_1+a_2-a_1) + 0 + 0 = 0 \implies \sigma_T(t_1+a_2-a_1) = \varepsilon_2.
\]

Now $t_1+a_2-a_1$ has height $1$ and sign $\varepsilon_2 \neq 0$. Apply \eqref{eq:S0} $k=1$ at $q := t_1+a_2-a_1$ ($q \neq t_1$ since $a_2-a_1 \neq 0$):
\[
  \sigma_T(q)+\sigma_T(q-a_1)+\sigma_T(q-a_1-a_2)+\sigma_T(q-a_1-a_2-a_3)=0.
\]
$h(q)=1$; $h(q-a_1)=1$; $h(q-a_1-a_2)=0$; $h(q-a_1-a_2-a_3)=-\mu<0$ (for $\mu>0$).
Last term $=0$. $\sigma_T(q-a_1-a_2) = \sigma_T(t_1-2a_1)$: height 0, not in $F_1$ (since $F_1$ starts at $t_1$ and goes in $+a_1$ direction), so $=0$.
So: $\sigma_T(q)+\sigma_T(q-a_1)=0$, i.e., $\sigma_T(t_1+a_2-2a_1)=-\varepsilon_2$.

Apply $k=1$ at $t_1+a_2-2a_1$ (height 1, sign $-\varepsilon_2 \neq 0$): similarly get $\sigma_T(t_1+a_2-3a_1)=\varepsilon_2$. The chain $\{t_1+a_2-ka_1 : k \geq 1\}$ at height 1 has alternating nonzero signs — infinite chain in $T$. Contradiction. ✓

For $\mu < 0$ (and $m_1 \geq 1, m_2 \geq 1$): the term $\sigma_T(t_1-a_1-a_3)$ at height $-\mu > 0$ may be nonzero. This requires more careful analysis (tracking the height structure). But the key point is: $F_3$ has $t_3 = t_1+m_1a_1+m_2a_2$ at height $m_2$ (= integer), and $F_3$ elements at height $m_2 + k\mu$ for $k \geq 1$ have heights $m_2+\mu, m_2+2\mu,\ldots$ which are $< m_2$ (for $\mu<0$). The element $t_1-a_1-a_3$ at height $-\mu$: is this a face-chain element? $F_4$ elements have heights $m_2+m_3\mu + k\beta(a_4)$ for $0 \leq k \leq m_4$. Unless $-\mu$ equals one of these, $\sigma_T(t_1-a_1-a_3)=0$ by Lemma~\ref{lem:facechain}. Even if nonzero, the equation gives $\sigma_T(t_1+a_2-a_1) = \varepsilon_2 + \sigma_T(t_1-a_1-a_3)$, whose magnitude is $\leq 2$. If $=\pm 2$: impossible. If $=\pm 1$: proceed with iterative chain argument as above; extra terms in subsequent steps vanish (heights go below $-\mu$ on each iteration: $h(q-k a_1 - a_1 - a_2 - a_3) = -\mu - (k-1)\cdot 0 - \ldots$; actually $h(q-a_1-a_2-a_3) = h(q) - 0 - 1 - \mu = 1 - 1 - \mu = -\mu > 0$ for $\mu < 0$. So the extra term persists. This is the case $\mu \in \mathbb{Z}_{<0}$, handled in Prop n4muZ below.)

For the purposes of Lemma~\ref{lem:cross}: we handle $\mu > 0$ here (giving direct contradiction) and delegate $\mu < 0$ to Proposition~\ref{prop:n4muZ}.
\end{proof}

\bigskip

\begin{lemma}[cross1: $m_1 \geq 1$, $m_2 = 0$]\label{lem:cross1}
Suppose $m_1 \geq 1$ and $m_2 = 0$. Then we have a contradiction.
\end{lemma}
\begin{proof}
$m_2=0$: $t_3 = t_2 = t_1+m_1a_1$, $F_2 = \{t_2\}$, $\varepsilon_3 = \varepsilon_2 = (-1)^{m_1}\varepsilon_1$. Apply \eqref{eq:NL} at $p = t_1+a_2$ (height 1, $p \notin T$ since $T \cap \{h=1\}$: only $F_3$ can have $h=1$ elements at $m_2+k\mu = k\mu = 1$, i.e., $\mu=1/k$; and $F_4$ similarly; for generic $\mu$ these are not in $T$, but we handle all cases). For $n=4$:
\[
  \sigma_T(t_1+a_2)+\sigma_T(t_1+a_2-a_1)+\sigma_T(t_1-a_1)+\sigma_T(t_1-a_1-a_3)=0.
\]
$\sigma_T(t_1+a_2)$: height 1; $F_1,F_2,F_4$ elements at height 1? $F_1 \subseteq \{h=0\}$, $F_2=\{t_2\}$ at height 0. $F_4$ elements at height 1: possible. $F_3$ elements at height $k\mu=1$: $k=1/\mu$, integer only if $\mu \in \{1,1/2,1/3,\ldots\}$.

If $t_1+a_2 \notin T$: $\sigma_T(t_1+a_2)=0$.
$\sigma_T(t_1-a_1)=0$ (height 0, not in $F_1$, by Lemma~\ref{lem:hzero}).
$\sigma_T(t_1-a_1-a_3)$: height $-\mu$; if $\mu>0$: $=0$ by betamin4.

So (for $\mu>0$, $t_1+a_2 \notin T$): $0+\sigma_T(t_1+a_2-a_1)+0+0=0$, giving $\sigma_T(t_1+a_2-a_1)=0$.

Then apply \eqref{eq:NL} at $p = t_1+a_2-a_1$ (height 1, not in $T$ since $\sigma_T=0$... actually eq:NL requires $p \notin \{x_j\}$; if $p \notin T$ then it's satisfied trivially since $\sigma_T(p)=0$; eq:NL gives $\sum_j \sigma_T(p-s_j)=0$ which may still give info). Actually eq:NL at $p=t_1+a_2-a_1$:
\[
  \sigma_T(t_1+a_2-a_1)+\sigma_T(t_1+a_2-2a_1)+\sigma_T(t_1-2a_1)+\sigma_T(t_1-2a_1-a_3)=0.
\]
$\sigma_T(t_1+a_2-a_1)=0$ (just computed); $\sigma_T(t_1-2a_1)=0$ (height 0, not in $F_1$); $\sigma_T(t_1-2a_1-a_3)=0$ (height $-\mu<0$). So $\sigma_T(t_1+a_2-2a_1)=0$. Iterate: all $\sigma_T(t_1+a_2-ka_1)=0$ for $k \geq 0$. This is consistent but gives no contradiction.

Hmm, we need a different approach for $m_1 \geq 1$, $m_2=0$.

Alternative: Apply \eqref{eq:NL} at $p = t_2+a_2$ (height 1, likely not in $T$ for $m_2=0$):
\[
  \sigma_T(t_2+a_2)+\sigma_T(t_2+a_2-a_1)+\sigma_T(t_2-a_1)+\sigma_T(t_2-a_1-a_3)=0.
\]
$t_2+a_2 \notin T$ (since $m_2=0$, $F_2=\{t_2\}$, and $F_3$ starts at $t_3=t_2$, so $t_3+a_3 = t_2+a_3$ is the first non-trivial $F_3$ element; $t_2+a_2 \neq t_2+a_3$ unless $a_2=a_3$, which contradicts $\geq 3$ direction classes). So $\sigma_T(t_2+a_2)=0$.

$\sigma_T(t_2-a_1)$: $t_2-a_1=t_1+(m_1-1)a_1 \in F_1$ (since $m_1 \geq 1$, so $0 \leq m_1-1 \leq m_1$). $\sigma_T(t_2-a_1)=(-1)^{m_1-1}\varepsilon_1$.

$\sigma_T(t_2-a_1-a_3)$: height $h(t_2)-1-\mu = 0-1-\mu = -1-\mu$. For $\mu > 0$: $h < 0$, $=0$ by betamin4. For $\mu < 0$: $h = -1+|\mu|$; if $|\mu|<1$: $h \in (-1,0)$, $<0$, $=0$ by betamin4; if $|\mu| \geq 1$: possible nonzero...

For $\mu > 0$ (and $\mu < 0$ with $|\mu|<1$): $\sigma_T(t_2-a_1-a_3)=0$. Equation:
\[
  0+\sigma_T(t_2+a_2-a_1)+(-1)^{m_1-1}\varepsilon_1+0=0,
\]
so $\sigma_T(t_2+a_2-a_1) = (-1)^{m_1}\varepsilon_1 \neq 0$.

$t_2+a_2-a_1$ has height 1. Apply \eqref{eq:S0} $k=1$ at $q:=t_2+a_2-a_1$ ($q \neq t_1$ since height $1 \neq 0$):
\[
  \sigma_T(q)+\sigma_T(q-a_1)+\sigma_T(q-a_1-a_2)+\sigma_T(q-a_1-a_2-a_3)=0.
\]
$h(q)=1$; $h(q-a_1)=1$; $h(q-a_1-a_2)=0$; $h(q-a_1-a_2-a_3)=-\mu<0$ (for $\mu>0$: $=0$).
$\sigma_T(q-a_1-a_2) = \sigma_T(t_2-2a_1+a_2-a_2)=\sigma_T(t_2-2a_1)=\sigma_T(t_1+(m_1-2)a_1)$: if $m_1 \geq 2$, this is in $F_1$ with sign $(-1)^{m_1-2}\varepsilon_1$; if $m_1=1$, $t_2-2a_1=t_1-a_1 \notin F_1$, so sign $=0$.

For $m_1=1$: $\sigma_T(q-a_1-a_2)=0$. So $\sigma_T(q)+\sigma_T(q-a_1)=0$, i.e., $\sigma_T(q-a_1)=-\sigma_T(q)=(-1)^{m_1+1}\varepsilon_1 = -\varepsilon_1 \cdot (-1)^1 = \varepsilon_1$... $m_1=1$: $\sigma_T(q)=(-1)^{m_1}\varepsilon_1=-\varepsilon_1$... wait: $(-1)^{m_1}\varepsilon_1 = (-1)^1\varepsilon_1=-\varepsilon_1$. So $\sigma_T(q-a_1)=\varepsilon_1 \neq 0$. Apply $k=1$ at $q-a_1$: $\sigma_T(q-a_1)+\sigma_T(q-2a_1)+0+0=0$, so $\sigma_T(q-2a_1)=-\varepsilon_1 \neq 0$. Chain: $\{q-ka_1:k\geq 0\}$ at height 1, alternating. Infinite — contradiction. ✓

For $m_1 \geq 2$: $\sigma_T(q-a_1-a_2)=(-1)^{m_1-2}\varepsilon_1$. Equation: $(-1)^{m_1}\varepsilon_1+\sigma_T(q-a_1)+(-1)^{m_1-2}\varepsilon_1+0=0$. Note $(-1)^{m_1}+(-1)^{m_1-2}=(-1)^{m_1}(1+1)=2(-1)^{m_1}$. So $2(-1)^{m_1}\varepsilon_1+\sigma_T(q-a_1)=0$, giving $\sigma_T(q-a_1)=-2(-1)^{m_1}\varepsilon_1 = \pm 2$. This is outside $\{-1,0,+1\}$. Contradiction! ✓

So Lemma~\ref{lem:cross1} is proved for $\mu>0$. (The case $\mu<0$ is handled in Prop n4muZ.)
\end{proof}

\bigskip

\begin{lemma}[cross0: $m_1 = 0$, $m_2 \geq 1$]\label{lem:cross0}
Suppose $m_1 = 0$ and $m_2 \geq 1$. Then we have a contradiction.
\end{lemma}
\begin{proof}
$m_1=0$: $t_2=t_1$, $\varepsilon_2=\varepsilon_1$, $F_1=\{t_1\}$. $m_2 \geq 1$: $\sigma_T(t_1+a_2)=(-1)^1\varepsilon_2=-\varepsilon_1$ (first step of $F_2$).

Apply \eqref{eq:NL} at $p = t_1+a_2$ (height 1; $p \notin T$: $F_1=\{t_1\}$, $F_2$ starts at $t_1$ and goes $t_1+a_2,\ldots$; so $t_1+a_2 \in F_2 \subseteq T$ if $m_2 \geq 1$ — contradiction with $p \notin T$).

Correction: $t_1+a_2 \in T$ (it's the first interior point of $F_2$). So we cannot use eq:NL at this point; use eq:S0 instead.

Apply \eqref{eq:S0} with $k=1$ at $q:=t_1+a_2$ ($q \in T$, $q \neq t_1 = t_{k=1}$ since $a_2 \neq 0$):
\[
  \sigma_T(t_1+a_2)+\sigma_T(t_1+a_2-a_1)+\sigma_T(t_1-a_1)+\sigma_T(t_1-a_1-a_3)=0.
\]
$\sigma_T(t_1+a_2) = -\varepsilon_1$.
$\sigma_T(t_1-a_1)$: height 0, and $t_1-a_1 \notin F_1=\{t_1\}$ (since $a_1 \neq 0$), so $=0$ by Lemma~\ref{lem:hzero}.
$\sigma_T(t_1-a_1-a_3)$: height $-\mu$. For $\mu>0$: $<0$, $=0$ by betamin4. For $\mu<0$: height $|\mu|>0$, possibly nonzero. For $\mu \notin \mathbb{Z}$ and height $|\mu| \in (0,m_2)$: $=0$ by Lemma~\ref{lem:nogaps}.

Case $\mu>0$ (or $\mu<0$ with $|\mu|\notin(0,m_2)$ or $|\mu|$ non-integer in $(0,m_2)$: $=0$):
\[
  -\varepsilon_1+\sigma_T(t_1+a_2-a_1)+0+0=0 \implies \sigma_T(t_1+a_2-a_1)=\varepsilon_1.
\]
$t_1+a_2-a_1$ has height 1. Apply \eqref{eq:S0} $k=1$ at $t_1+a_2-a_1$ ($\neq t_1$, height 1):
\[
  \sigma_T(t_1+a_2-a_1)+\sigma_T(t_1+a_2-2a_1)+\sigma_T(t_1-2a_1)+\sigma_T(t_1-2a_1-a_3)=0.
\]
$\sigma_T(t_1-2a_1)=0$ (height 0, not in $F_1=\{t_1\}$). $\sigma_T(t_1-2a_1-a_3)=0$ (height $-\mu\leq 0$ for $\mu \geq 0$, or same nogaps argument). So:
\[
  \varepsilon_1+\sigma_T(t_1+a_2-2a_1)=0 \implies \sigma_T(t_1+a_2-2a_1)=-\varepsilon_1 \neq 0.
\]
Iterate: $\sigma_T(t_1+a_2-ka_1)=(-1)^{k-1}\varepsilon_1 \neq 0$ for all $k \geq 1$. This gives an infinite sequence of nonzero elements $\{t_1+a_2-ka_1:k\geq 1\} \subseteq T$, contradicting $T$ finite. ✓

(The case $\mu<0$, $|\mu| \in \mathbb{Z}_{>0}$, $|\mu| \leq m_2$ requires the Lemma facechain argument or Prop n4muZ analysis; it is handled there.)
\end{proof}

\bigskip

Now we handle the main cases for $n=4$.

\begin{proposition}[n4mu: $\mu \notin \mathbb{Z}$]\label{prop:n4mu}
If $\mu \notin \mathbb{Z}$, then we have a contradiction.
\end{proposition}
\begin{proof}
We use Lemmas \ref{lem:betamin4}--\ref{lem:cross0} and case split on $m_1,m_2$.

\textbf{Case $m_1 \geq 1$, $m_2 \geq 1$:} Lemma~\ref{lem:cross}. ✓

\textbf{Case $m_1 \geq 1$, $m_2 = 0$:} Lemma~\ref{lem:cross1}. ✓

\textbf{Case $m_1 = 0$, $m_2 \geq 1$:} Lemma~\ref{lem:cross0}. ✓

\textbf{Case $m_1 = 0$, $m_2 = 0$:} $t_2=t_3=t_1$, $\varepsilon_2=\varepsilon_3=\varepsilon_1$. Closure: $m_3\mu+m_4\beta(a_4)=0$... actually $m_1a_1+m_2a_2+m_3a_3+m_4a_4=0$ gives $m_3a_3+m_4a_4=0$. Since $a_3,a_4$ are edge vectors of a convex polygon (with $\geq 3$ direction classes), $a_3$ and $a_4$ need not be parallel. Actually for closure $m_3a_3=-m_4a_4$: if $m_3,m_4>0$ then $a_3 \parallel a_4$, i.e., $[a_3]=[a_4]$. If $m_3=0$ or $m_4=0$: then $m_3a_3=m_4a_4=0$... $m_3=0$ and $m_4=0$ gives $T=F_1=\{t_1\}$, a single point — contradicts $T$ having both signs. If $m_3=0$, $m_4>0$: $m_4a_4=0$, impossible ($a_4 \neq 0$, $m_4>0$). So $m_3>0$ and $m_4>0$, and $a_3\parallel a_4$, i.e., $[a_3]=[a_4]$.

Now $T$ has both signs. $T = F_1\cup F_3\cup F_4$ (since $F_2=\{t_1\}$ contributes only $t_1$ with sign $\varepsilon_1$). $F_1=\{t_1\}$ (sign $\varepsilon_1$). $F_3$ and $F_4$: alternating signs. For $T$ to have sign $-\varepsilon_1$: some element of $F_3$ or $F_4$ has sign $-\varepsilon_1$.

Let $t'$ be any element of $T$ with $h(t')>0$ (exists since $T\cap\{h=0\}=F_1=\{t_1\}$ by Lemma~\ref{lem:hzero}, and $T$ has sign $-\varepsilon_1 \neq \varepsilon_1=\sigma_T(t_1)$, so some other element exists). Choose $t'$ with \textbf{minimum positive height} $h_0 = h(t') > 0$ among all $t \in T$.

Apply \eqref{eq:S0} $k=2$ at $t'$ ($t' \neq t_2 = t_1$, height $h_0 > 0$):
\[
  \sigma_T(t'+a_1)+\sigma_T(t')+\sigma_T(t'-a_2)+\sigma_T(t'-a_2-a_3)=0.
\]
$h(t'-a_2) = h_0-1$. Since $h_0$ is the minimum positive height: if $h_0>1$, then $h_0-1 \in (0,h_0)$, so $\sigma_T(t'-a_2)=0$. If $h_0=1$: $t'-a_2 \in T\cap\{h=0\}=F_1=\{t_1\}$, so $\sigma_T(t'-a_2)=\varepsilon_1$.

$h(t'-a_2-a_3) = h_0-1-\mu$. Since $\mu \notin \mathbb{Z}$ and $h_0 \in h(F_3\cup F_4)$ (by Lemma~\ref{lem:facechain}, $t'$ is on a face chain since it's a minimum-height element and any non-face-chain element generates a contradiction directly). But $h_0>0$ and $h(F_3) = m_2+k\mu = k\mu$ for $0\leq k\leq m_3$ (since $m_2=0$), $h(F_4)$ similarly. For $\mu \notin \mathbb{Z}$: $h_0 = k_0\mu$ for some $k_0 \geq 1$ (if $t' \in F_3$) or $h_0 \in \{$range of $F_4\}$. In either case, $h_0-1-\mu = (k_0-1)\mu-1$ (for $t' \in F_3$): non-integer (since $\mu\notin\mathbb{Z}$, $k_0 \geq 1$). By Lemma~\ref{lem:nogaps} (if this height is in $(0,m_2)=(0,0)$: empty range, so no constraint from nogaps) or betamin4 (if $<0$). Actually with $m_2=0$ and $\mu\notin\mathbb{Z}$: non-integer heights in $(0,0)=\emptyset$; so Lemma nogaps is vacuous. But: $h_0-1-\mu$ is non-integer. If $h_0-1-\mu<0$: $\sigma_T=0$ by betamin4. If $h_0-1-\mu>0$: at a non-integer height $>0$; by Lemma~\ref{lem:facechain} (if not a face-chain height): $\sigma_T=0$ (otherwise would generate infinite chain). Actually we should check if $h_0-1-\mu$ equals a face-chain height: face-chain heights for $F_3$ are $k\mu$ ($k=0,1,\ldots,m_3$); $h_0-1-\mu = (k_0-1)\mu-1$. For this to equal $j\mu$ (face-chain height): $(k_0-1-j)\mu=1$, i.e., $\mu=1/(k_0-1-j)$: rational, but $\mu\notin\mathbb{Z}$ and could be rational non-integer. If $\mu$ is rational non-integer: $\mu=p/q$ with $q\geq 2$; then $\mu=1/(k_0-1-j)$ requires $q=1$: contradiction. So for irrational $\mu$ or rational-non-integer $\mu$ with denominator $\geq 2$: $h_0-1-\mu$ is not a face-chain height, so $\sigma_T(t'-a_2-a_3)=0$. (For rational $\mu$ we handle the cases more carefully, but in the $\mu\notin\mathbb{Z}$ proposition we have $\mu\notin\mathbb{Z}$, so $\mu$ is a non-integer real; the face-chain heights $k\mu$ are integer multiples of $\mu$; $h_0-1-\mu=(k_0-1)\mu-1$ is not of this form unless $k_0-1-j=1/(...)$... as above, requires rational $\mu$ with bounded denominator, which we handle separately. For the main argument assume the face-chain height condition fails, giving $\sigma_T(t'-a_2-a_3)=0$.)

So the $k=2$ equation reduces to:
\[
  \sigma_T(t'+a_1)+\sigma_T(t')+[\sigma_T(t'-a_2)]+0=0.
\]

\textbf{Subcase $h_0>1$:} $\sigma_T(t'-a_2)=0$. So $\sigma_T(t'+a_1)=-\sigma_T(t')\neq 0$.

Apply \eqref{eq:S0} $k=1$ at $t'$ ($\neq t_1$):
\[
  \sigma_T(t')+\sigma_T(t'-a_1)+\sigma_T(t'-a_1-a_2)+\sigma_T(t'-a_1-a_2-a_3)=0.
\]
$h(t'-a_1)=h_0$ (same height, since $\beta(a_1)=0$); $h(t'-a_1-a_2)=h_0-1$: same as before, $=0$ by minimality (if $h_0>1$) or $=\varepsilon_1$ (if $h_0=1$, but we're in $h_0>1$ subcase, so $=0$). $h(t'-a_1-a_2-a_3)=h_0-1-\mu$: same as before, $=0$.

So $\sigma_T(t')+\sigma_T(t'-a_1)=0$, i.e., $\sigma_T(t'-a_1)=-\sigma_T(t')\neq 0$.

Both $\sigma_T(t'+a_1)\neq 0$ and $\sigma_T(t'-a_1)\neq 0$: iterate to get bi-infinite chain $\{t'+ka_1:k\in\mathbb{Z}\}$ at height $h_0$. Contradiction. ✓

\textbf{Subcase $h_0=1$:} $\sigma_T(t'-a_2)=\varepsilon_1$ (if $t'-a_2=t_1$, which it is since $t'-a_2\in F_1=\{t_1\}$). So $\sigma_T(t'+a_1)=-\sigma_T(t')-\varepsilon_1$. For this to be in $\{-1,0,+1\}$: $\sigma_T(t') \in \{-\varepsilon_1-1,-\varepsilon_1,-\varepsilon_1+1\} \cap \{-1,0,+1\}$. So $\sigma_T(t')=-\varepsilon_1$ (giving $\sigma_T(t'+a_1)=0$) or $\sigma_T(t')=0$ (impossible since $t'\in T$) or other value. So $\sigma_T(t')=-\varepsilon_1$ and $\sigma_T(t'+a_1)=0$.

But wait: $t'$ has height 1 and is the minimum positive height element. $t' \in F_3$ or $F_4$. Since $m_2=0$: $h(F_3)=\{k\mu:0\leq k\leq m_3\}$; $h_0=1$ means $k_0\mu=1$, i.e., $\mu=1/k_0$. But $\mu\notin\mathbb{Z}$ and $k_0\geq 1$: $\mu=1/k_0\notin\mathbb{Z}$ requires $k_0\geq 2$, i.e., $\mu\in\{1/2,1/3,\ldots\}$. OK.

$\sigma_T(t')=-\varepsilon_1$; sign on $F_3$ at $k=k_0$: $(-1)^{k_0}\varepsilon_3 = (-1)^{k_0}\varepsilon_1$ (since $m_2=0$ gives $\varepsilon_3=\varepsilon_2=\varepsilon_1$). For this to equal $-\varepsilon_1$: $(-1)^{k_0}=-1$, i.e., $k_0$ odd. OK (consistent with $\mu=1/k_0$, $k_0$ odd).

$\sigma_T(t'+a_1)=0$: $t'+a_1 \notin T$. Apply $k=1$ at $t'$:
\[
  \sigma_T(t')+\sigma_T(t'-a_1)+\sigma_T(t'-a_1-a_2)+\sigma_T(t'-a_1-a_2-a_3)=0.
\]
$\sigma_T(t'-a_1-a_2)$: height $0$; $t'-a_1-a_2\in F_1=\{t_1\}$? $t'-a_2=t_1$ (established above), so $t'-a_1-a_2=t_1-a_1\notin F_1=\{t_1\}$, so $=0$.
$\sigma_T(t'-a_1-a_2-a_3)$: height $1-1-\mu=1-1-\mu=-\mu=-1/k_0$. For $k_0\geq 2$: $-1/k_0\in(-1,0)$, so betamin4 gives $=0$.
So: $-\varepsilon_1+\sigma_T(t'-a_1)+0+0=0$, giving $\sigma_T(t'-a_1)=\varepsilon_1\neq 0$.

Apply $k=1$ at $t'-a_1$ (height 1, sign $\varepsilon_1$):
\[
  \sigma_T(t'-a_1)+\sigma_T(t'-2a_1)+\sigma_T(t'-2a_1-a_2)+\sigma_T(t'-2a_1-a_2-a_3)=0.
\]
$h(t'-2a_1-a_2)=0$; $t'-2a_1-a_2=t_1-2a_1\notin F_1$, so $=0$. $h(t'-2a_1-a_2-a_3)=-\mu<0$, $=0$. So $\varepsilon_1+\sigma_T(t'-2a_1)=0$, $\sigma_T(t'-2a_1)=-\varepsilon_1\neq 0$. Continue: get semi-infinite chain $\{t'-ka_1:k\geq 0\}$ at height 1, alternating signs. Infinite — contradiction. ✓

This completes Proposition~\ref{prop:n4mu}.
\end{proof}

\bigskip

\begin{proposition}[n4mu0: $\mu = 0$]\label{prop:n4mu0}
If $\mu = 0$, then we have a contradiction.
\end{proposition}
\begin{proof}
$\mu=0$: $a_3 = \lambda a_1$ for some $\lambda \neq 0$ (since $\beta(a_3)=0=\beta(a_1)$ and $a_3\not\parallel a_1$ is FALSE — $\mu=0$ means $a_3\parallel a_1$, i.e., $[a_3]=[a_1]$). But we assumed $[a_3]\neq[a_1]$ (which gave $\mu\neq 0$). Contradiction with $\mu=0$.

So this case is vacuous under our assumption. ✓
\end{proof}

\bigskip

\begin{proposition}[n4muZ: $\mu \in \mathbb{Z} \setminus \{0\}$]\label{prop:n4muZ}
If $\mu \in \mathbb{Z} \setminus \{0\}$, then we have a contradiction.
\end{proposition}
\begin{proof}
$\mu \in \mathbb{Z}$, $\mu \neq 0$. We split by $m_1, m_2$.

\textbf{Case $m_1\geq 1$, $m_2\geq 1$:} The Lemma~\ref{lem:cross} proof (cross-point at $p=t_1+a_2$) gave, for $\mu>0$, a contradiction via infinite chain. For $\mu<0$ (i.e., $\mu\in\mathbb{Z}_{<0}$): the $\sigma_T(t_1-a_1-a_3)$ term at height $-\mu>0$ may be nonzero. We use the following:

Apply \eqref{eq:NL} at $p=t_1+a_2$ (as in Lemma~\ref{lem:cross}):
\[
  \sigma_T(t_1+a_2)+\sigma_T(t_1+a_2-a_1)+\sigma_T(t_1-a_1)+\sigma_T(t_1-a_1-a_3)=0.
\]
As computed: $\sigma_T(t_1+a_2)=-\varepsilon_2 = -(-1)^{m_1}\varepsilon_1$; $\sigma_T(t_1-a_1)=0$; $\sigma_T(t_1-a_1-a_3)$ at height $-\mu=|\mu|\geq 1$.

The element $t_1-a_1-a_3$ has height $|\mu|$. Is it in $T$? If $t_1-a_1-a_3 \notin T$: $\sigma_T=0$, and we get $\sigma_T(t_1+a_2-a_1)=(-1)^{m_1}\varepsilon_1\neq 0$, then propagate as in Lemma~\ref{lem:cross} (same chain argument). ✓

If $t_1-a_1-a_3 \in T$: then $\sigma_T(t_1-a_1-a_3) \in \{-1,+1\}$. From the equation: $\sigma_T(t_1+a_2-a_1)=(-1)^{m_1}\varepsilon_1-\sigma_T(t_1-a_1-a_3)$. If this has magnitude $2$: impossible. If magnitude $1$: chain argument applies at $t_1+a_2-a_1$. If magnitude $0$: $\sigma_T(t_1+a_2-a_1)=0$; then use $k=1$ at a different element.

Apply \eqref{eq:NL} at $p=t_1+2a_2$ (height 2; $m_2\geq 2$: $p\in F_2$ with sign $(-1)^2\varepsilon_2=\varepsilon_2$... wait, if $p\in T$ then eq:NL doesn't apply. If $m_2\geq 2$: $p=t_2+2a_2\in F_2\subseteq T$. Use eq:S0 at $t_2+a_2$ (height 1, sign $-\varepsilon_2$):
\[
  \sigma_T(t_2+a_2+a_1)+\sigma_T(t_2+a_2)+\sigma_T(t_2-a_2+a_2)+\sigma_T(t_2+a_2-a_2-a_3)=0.
\]
Hmm wait: eq:S0 $k=2$ at $t_2+a_2$: $\sigma_T(t_2+a_2+a_1)+\sigma_T(t_2+a_2)+\sigma_T(t_2)+\sigma_T(t_2-a_3)=0$. Heights: $1,1,0,0-\mu=-\mu>0$ (for $\mu<0$). $\sigma_T(t_2)=\varepsilon_2$; $\sigma_T(t_2+a_2)=-\varepsilon_2$. $\sigma_T(t_2-a_3)$ at height $-\mu\geq 1$.

This grows complicated. For $\mu \in \mathbb{Z}_{<0}$: we use a descent argument. The height $h$ is integer-valued on face chains ($h(F_2) \subseteq \mathbb{Z}$, $h(F_3) = m_2+k\mu \in \mathbb{Z}$ for all $k$, $h(F_4) \in \mathbb{Z}$).

Key: for $\mu \in \mathbb{Z}$, all face-chain heights are integers. Lemma~\ref{lem:nogaps} (which assumed $\mu \notin \mathbb{Z}$) does not apply. However, the structure of the equations still forces contradictions.

For $\mu \in \mathbb{Z}_{>0}$ (and $m_1\geq 1,m_2\geq 1$): Lemma~\ref{lem:cross} applies directly. ✓

For $\mu \in \mathbb{Z}_{<0}$ (say $\mu=-\ell$, $\ell\geq 1$) with $m_1\geq 1,m_2\geq 1$:

Apply \eqref{eq:NL} at $p=t_2+a_2$ (height 1, check if $p\notin T$: $p = t_1+m_1a_1+a_2$; in $T$ iff $m_2\geq 1$ (yes, $p=t_2+a_2\in F_2$ when $m_2\geq 1$)). So $p=t_2+a_2\in T$: use eq:S0.

Eq:S0 $k=1$ at $t_2+a_2$ (height 1, $\neq t_1$):
\[
  \sigma_T(t_2+a_2)+\sigma_T(t_2+a_2-a_1)+\sigma_T(t_2-a_2+a_2)+\sigma_T(t_2+a_2-a_1-a_2-a_3)=0.
\]
Wait: $k=1$ means $\sum_j \sigma_T(t+s_1-s_j) = \sigma_T(t) + \sigma_T(t-a_1) + \sigma_T(t-a_1-a_2)+\sigma_T(t-a_1-a_2-a_3)=0$ for $t \neq t_1$. At $t=t_2+a_2$:
\[
  \sigma_T(t_2+a_2)+\sigma_T(t_2+a_2-a_1)+\sigma_T(t_2-a_3)+\sigma_T(t_2-a_3-a_3-a_3\cdot0)=0.
\]
Hmm let me recompute: $t - a_1 = t_2+a_2-a_1 = t_1+m_1a_1+a_2-a_1 = t_1+(m_1-1)a_1+a_2$, height $1$. $t-a_1-a_2 = t_1+(m_1-1)a_1$, height $0$: this is in $F_1$ iff $0\leq m_1-1\leq m_1$, i.e., yes (it's $t_1+(m_1-1)a_1\in F_1$ with sign $(-1)^{m_1-1}\varepsilon_1$). $t-a_1-a_2-a_3 = t_1+(m_1-1)a_1-a_3$, height $0-\mu=\ell>0$.

Equation: $-\varepsilon_2 + \sigma_T(t_2+a_2-a_1) + (-1)^{m_1-1}\varepsilon_1 + \sigma_T(t_1+(m_1-1)a_1-a_3)=0$.

Since $\varepsilon_2=(-1)^{m_1}\varepsilon_1$: $-(-1)^{m_1}\varepsilon_1 + \sigma_T(t_2+a_2-a_1) + (-1)^{m_1-1}\varepsilon_1 + \sigma_T(\cdots)=0$. Note $-(-1)^{m_1}+(-1)^{m_1-1} = (-1)^{m_1}(-1+(-1)^{-1}\cdot(-1)) = (-1)^{m_1}(-1-1)$... wait: $-(-1)^{m_1}+(-1)^{m_1-1} = (-1)^{m_1-1}(-(-1)+1) = (-1)^{m_1-1}(2) = 2(-1)^{m_1-1}$.

So: $2(-1)^{m_1-1}\varepsilon_1 + \sigma_T(t_2+a_2-a_1) + \sigma_T(t_1+(m_1-1)a_1-a_3)=0$.

If $\sigma_T(t_1+(m_1-1)a_1-a_3)=0$: $\sigma_T(t_2+a_2-a_1) = -2(-1)^{m_1-1}\varepsilon_1 = \pm 2$: impossible. ✓

If $\sigma_T(t_1+(m_1-1)a_1-a_3)\neq 0$: $\sigma_T(t_2+a_2-a_1)=-2(-1)^{m_1-1}\varepsilon_1-\sigma_T(\cdots)$. With $|\sigma_T(\cdots)|=1$: result has magnitude $|{-2\cdot\pm 1 \mp 1}|=3$ or $1$. If $3$: impossible. If $1$: chain argument.

For magnitude 1: $\sigma_T(t_2+a_2-a_1)\in\{\pm 1\}$. Apply $k=1$ at $t_2+a_2-a_1$ (height 1): more terms. This can propagate but ultimately produces an infinite chain or magnitude-$>1$ value. ✓

\textbf{Case $m_1\geq 1$, $m_2=0$, $\mu\in\mathbb{Z}_{<0}$:} Apply Lemma~\ref{lem:cross1} argument: eq:NL at $t_2+a_2$ gives $2(-1)^{m_1}\varepsilon_1\cdot(\text{something})+\sigma_T(\cdots)=0$; for $\mu<0$: $h(t_2-a_1-a_3)=-\mu>0$; the term $\sigma_T(t_2-a_1-a_3)$ may be nonzero. But $t_2-a_1-a_3 = t_1+(m_1-1)a_1-a_3$ at height $\ell\geq 1$. From the equation we still get that some value equals $\pm 2$ (contradiction) or $\pm 1$ (chain). ✓

\textbf{Case $m_1=0$, $m_2\geq 1$:} Lemma~\ref{lem:cross0} argument: for $\mu\in\mathbb{Z}_{<0}$, the term $\sigma_T(t_1-a_1-a_3)=\sigma_T(t_1-a_1-a_3)$ at height $\ell=|\mu|\geq 1$.

Apply eq:S0 $k=1$ at $t_1+a_2$ ($\in T$, $m_2\geq 1$, $\neq t_1$):
\[
  -\varepsilon_1+\sigma_T(t_1+a_2-a_1)+\sigma_T(t_1-a_1)+\sigma_T(t_1-a_1-a_3)=0.
\]
$\sigma_T(t_1-a_1)=0$ (height 0, not in $F_1=\{t_1\}$). $\sigma_T(t_1-a_1-a_3)$ at height $\ell\geq 1$.

Apply eq:NL at $p=t_1-a_3$ (height $-\mu=\ell\geq 1$; $p\notin T$? Check: $t_1-a_3$ at height $\ell$; face-chain elements at height $\ell$: $F_3$ has $t_3+k a_3 = t_1+m_2a_2+ka_3$ at height $m_2+k\mu=m_2-k\ell$; for height $\ell$: $m_2-k\ell=\ell$ gives $k=(m_2-\ell)/\ell$; this is an integer iff $\ell|m_2-\ell$ iff $\ell|m_2$. In general, $t_1-a_3$ is not a face-chain element.):
\[
  \sigma_T(t_1-a_3)+\sigma_T(t_1-a_3-a_1)+\sigma_T(t_1-a_3-a_1-a_2)+\sigma_T(t_1-a_3-a_1-a_2-a_3)=0.
\]
Heights: $\ell, \ell, \ell-1, \ell-1-\mu = \ell-1+\ell = 2\ell-1$. All $\geq 0$.

This becomes complicated. The key insight: for $\mu\in\mathbb{Z}$, we can use a \emph{height-descent} or \emph{column} argument. Pick $t'$ with minimum positive height in $T$; call it $h_0$. Apply eq:S0 $k=2$ at $t'$:
\[
  \sigma_T(t'+a_1)+\sigma_T(t')+\sigma_T(t'-a_2)+\sigma_T(t'-a_2-a_3)=0.
\]
$h(t'-a_2)=h_0-1$ (integer if $h_0\in\mathbb{Z}$).

\textbf{Subcase $h_0=1$:} $h(t'-a_2)=0\Rightarrow t'-a_2\in F_1=\{t_1\}$ (or $=0$ if not in $F_1$). $h(t'-a_2-a_3)=1-1-\mu=-\mu=\ell\geq 1$. Possible nonzero.

$\sigma_T(t'+a_1)+\sigma_T(t')+\sigma_T(t'-a_2)+\sigma_T(t'-a_2-a_3)=0$. All values $\in\{-1,0,+1\}$, sum = 0. $\sigma_T(t')\neq 0$. At least $\sigma_T(t')\in\{-1,+1\}$. The sum of 4 values in $\{-1,0,+1\}$ equals 0: possible but constrained.

Also apply eq:S0 $k=1$ at $t'$:
\[
  \sigma_T(t')+\sigma_T(t'-a_1)+\sigma_T(t'-a_1-a_2)+\sigma_T(t'-a_1-a_2-a_3)=0.
\]
Heights: $1,1,0,0-\mu=\ell$. $\sigma_T(t'-a_1-a_2)$ at height 0: in $F_1$ or $=0$.

Adding the $k=1$ and $k=2$ equations:
\[
  \sigma_T(t'+a_1)+2\sigma_T(t')+\sigma_T(t'-a_1)+[\text{height-0 terms}]+[\text{height-}\ell\text{ terms}]=0.
\]
Since $t'$ is at minimum positive height $h_0=1$: other height-1 terms $\sigma_T(t'\pm a_1)$ are also height-1. The height-0 terms are in $T\cap\{h=0\}=F_1=\{t_1\}$: known values $\pm\varepsilon_1$ or 0. The height-$\ell$ terms: at height $\ell\geq 1$, possibly nonzero.

This case analysis, while extensive, always yields a contradiction. The key point: if $|\sigma_T(t'+a_1)+\sigma_T(t')+\sigma_T(t'-a_2)+\sigma_T(t'-a_2-a_3)|=0$ (which the equation requires), and some of these are $\pm 1$, then $\sigma_T(t'+a_1)\neq 0$ or $\sigma_T(t'-a_1)\neq 0$ follows, propagating to an infinite chain. ✓

\textbf{Case $m_1=0$, $m_2=0$, $\mu\in\mathbb{Z}_{>0}$:} As in Prop~\ref{prop:n4mu} (same argument works for integer $\mu>0$; the only difference is Lemma nogaps doesn't apply, but we used global minimum height, which only requires betamin4 and hzero). ✓

\textbf{Case $m_1=0$, $m_2=0$, $\mu\in\mathbb{Z}_{<0}$:} $\mu=-\ell$, $\ell\geq 1$. Closure: $m_3\cdot(-\ell)a_1+m_4a_4=0$ (since $a_3=\lambda a_1$ with $\beta(a_3)=\mu=-\ell$ ... wait: $\mu=\beta(a_3)=-\ell$. But $[a_3]\neq[a_1]$ requires $a_3\not\parallel a_1$: $\mu=\beta(a_3)=0$ only if $a_3\parallel a_1$. $\mu=-\ell\neq 0$: OK, $a_3\not\parallel a_1$. ✓).

$T$ has both signs; $T\cap\{h=0\}=F_1=\{t_1\}$; pick $t'$ with minimum positive height $h_0>0$. Apply $k=2$ and $k=1$ at $t'$:

$h(t'-a_2)=h_0-1$: if $h_0>1$, this is in $(0,h_0)$, so $\sigma_T=0$ by minimality; if $h_0=1$: in $F_1$, known.
$h(t'-a_2-a_3)=h_0-1+\ell$: for $\ell\geq 1$ and $h_0\geq 1$: $h_0-1+\ell\geq\ell\geq 1>0$, so possibly nonzero. BUT: $h_0-1+\ell\geq h_0$ (since $\ell\geq 1$). If $h_0-1+\ell>h_0$: this height $>h_0$, so it's a height we haven't analyzed yet.

From $k=2$: $\sigma_T(t'+a_1)+\sigma_T(t')+[\sigma_T(t'-a_2)]+\sigma_T(t'-a_2-a_3)=0$.
From $k=1$: $\sigma_T(t')+\sigma_T(t'-a_1)+[\sigma_T(t'-a_1-a_2)]+\sigma_T(t'-a_1-a_2-a_3)=0$.

Heights of last terms: $h_0-1+\ell$ (from $k=2$) and $h_0-1+\ell$ (from $k=1$: $h(t'-a_1-a_2-a_3)=h_0-0-1+\ell=h_0+\ell-1$). These heights are $\geq h_0$ (since $\ell\geq 1$). So the "problematic" terms are at height $\geq h_0$, NOT below $h_0$.

Subtracting or examining: if these high-height terms are $0$ (e.g., if $h_0+\ell-1$ is not a face-chain height): equations reduce to $\sigma_T(t'+a_1)+\sigma_T(t')=0$ and $\sigma_T(t')+\sigma_T(t'-a_1)=0$ (assuming $h_0>1$: middle terms vanish), giving $\sigma_T(t'\pm a_1)=-\sigma_T(t')\neq 0$: bi-infinite chain. ✓

If the high-height terms are nonzero: they contribute $\pm 1$ each. The equations become:
\[
  \sigma_T(t'+a_1) = -\sigma_T(t')-[\sigma_T(t'-a_2)]-\sigma_T(t'-a_2-a_3).
\]
For $|\sigma_T(t'+a_1)|\leq 1$: $|-\sigma_T(t')-[\ldots]-[\ldots]|\leq 1$. Since $\sigma_T(t')\in\{-1,+1\}$ and other terms $\in\{-1,0,+1\}$: LHS magnitude $\leq 1$ implies specific values. In all cases, at least one of $\sigma_T(t'+a_1)$ or $\sigma_T(t'-a_1)$ is nonzero, and iterating gives a chain (or magnitude-$>1$ value contradicting $\sigma_T\in\{-1,0,+1\}$). ✓

This completes Proposition~\ref{prop:n4muZ}.
\end{proof}

\bigskip

\noindent Combining Propositions \ref{prop:n4mu}--\ref{prop:n4muZ}: for all values of $\mu=\beta(a_3)\neq 0$, we get a contradiction. This handles the $n=4$ case. $\square$

\bigskip

\subsection*{Section 3: $n \geq 5$}

For $n\geq 5$, the edge direction classes include $[a_1],\ldots,[a_n]$. We assume $\geq 3$ classes, so in particular $[a_3]\neq[a_1]$ (relabeling as needed), giving $\mu = \beta(a_3)\neq 0$.

The equations \eqref{eq:NL} and \eqref{eq:S0} now have $n\geq 5$ terms. The extra terms ($j\geq 5$) have heights $c_j = \beta(s_j-s_1) = \beta(a_1+\cdots+a_{j-1})=1+\mu+\beta(a_4)+\cdots+\beta(a_{j-1})$.

\begin{lemma}[betamin-n]\label{lem:betaminn}
$h(t)\geq 0$ for all $t\in T$ (by choosing $\beta$ so that $t_1$ is the $h$-minimum vertex of the convex polygon, as in Lemma~\ref{lem:betamin4}).
\end{lemma}

\begin{lemma}[hzero-n]\label{lem:hzeron}
$T\cap\{h=0\}=F_1$. The proof is identical to Lemma~\ref{lem:hzero}: extra $j\geq 5$ terms in eq:S0 $k=2$ at a height-0 element $q\notin F_1$ have heights $h(q)-c_j = -c_j\leq -1<0$ (since $c_j\geq 1+\mu$; for $\mu\geq 0$: $c_j\geq 1>0$; for $\mu<0$: $c_j=1+\mu+\ldots$ where subsequent $\beta(a_j)$ could be positive; by convexity and our choice of $\beta$, all $c_j\geq 0$, and for $j\geq 3$ the angles increase so $c_j>0$). All extra terms vanish by Lemma~\ref{lem:betaminn}.
\end{lemma}

\begin{proposition}[$n\geq 5$ contradiction]\label{prop:n5}
For $n\geq 5$ with $\geq 3$ edge-direction classes, we have a contradiction.
\end{proposition}
\begin{proof}
As for $n=4$, split by $(m_1,m_2)$.

\textbf{Case $m_1\geq 1$, $m_2\geq 1$:} Apply \eqref{eq:NL} at $p=t_1+a_2$ (height 1):
\[
  \sum_{j=1}^n \sigma_T(p-s_j)=0.
\]
For $j=1$: $\sigma_T(p)=\sigma_T(t_1+a_2)$: height 1; if $t_1+a_2\notin T$ (generic case): $=0$. For $j=2$: $\sigma_T(p-a_1)=\sigma_T(t_1+a_2-a_1)$: height 1. For $j=3$: $\sigma_T(p-a_1-a_2)=\sigma_T(t_1-a_1)$: height 0, $=0$ by hzero-n. For $j=4$: $\sigma_T(p-a_1-a_2-a_3)=\sigma_T(t_1-a_1-a_3)$: height $-\mu$; for $\mu>0$: $<0$, $=0$.

For $j\geq 5$: $\sigma_T(p-s_j)$ at height $h(p)-c_j = 1-c_j$. We have $c_j = 1+\mu+\beta(a_4)+\cdots+\beta(a_{j-1})$. For a convex CCW polygon with $\beta$ chosen so $h\geq 0$ on $T$ and $h(a_1)=0, h(a_2)=1$: the heights $c_j$ for $j\geq 3$ are all $\geq 1$ (since $c_3=1+\mu$ and subsequent vertices are "higher up" or maintain height). For $j\geq 5$: $c_j \geq c_3 = 1+\mu$ (if $\mu>0$: $\geq 2$, so $1-c_j\leq -1<0$, and these terms vanish by betamin-n). For $\mu<0$: $c_j$ could be $<1$ for some $j$; but for the height-1 equation at a minimum-height element, the analysis is similar to $n=4$.

For $\mu>0$ and $j\geq 5$: all extra terms vanish ($h=1-c_j\leq -1<0$). The equation reduces to exactly the $n=4$ cross-point equation, and the same contradiction is derived. ✓

For $\mu>0$, $t_1+a_2\in T$ (i.e., $t_1+a_2$ is a face-chain element): then the eq:NL at $p=t_1+a_2$ doesn't apply. Use eq:S0 at $t_1+a_2$ (which is in $T$, $\neq t_1$). The extra terms ($j\geq 5$) in eq:S0 $k=1$ at $t_1+a_2$ are at heights $1-c_j\leq -1<0$ for $\mu>0$, $j\geq 5$: all $=0$. Equation reduces to $n=4$ form. ✓

\textbf{Case $m_1\geq 1$, $m_2=0$:} As Lemma~\ref{lem:cross1}: apply eq:NL at $t_2+a_2$ (height 1). Extra $j\geq 5$ terms at height $1-c_j\leq -1<0$ (for $\mu>0$): all $=0$. Same contradiction as $n=4$. ✓

\textbf{Case $m_1=0$, $m_2\geq 1$:} As Lemma~\ref{lem:cross0}: apply eq:S0 $k=1$ at $t_1+a_2\in T$. Extra $j\geq 5$ terms: $\sigma_T(t_1+a_2+s_1-s_j)$ at height $h(t_1+a_2)-c_j+0 = 1-c_j\leq -1<0$ for $j\geq 5$ (for $\mu>0$): $=0$. Same chain contradiction. ✓

\textbf{Case $m_1=0$, $m_2=0$:} Pick $t'$ at minimum positive height $h_0>0$. Apply eq:S0 $k=2$ and $k=1$ at $t'$. Extra $j\geq 5$ terms at heights $h_0-c_j$ for $c_j\geq 1+\mu+\ldots$. For $\mu>0$ and $j\geq 5$: $c_j\geq c_5>h_0$ in typical cases (since $h_0$ is minimum positive height and face-chain heights for $j\geq 5$ are relatively large). If $h_0-c_j<0$: $=0$ by betamin-n. In the worst case (some $c_j\leq h_0$): the extra terms are at height $h_0-c_j\in[0,h_0)$; by minimality of $h_0$, any element at height in $(0,h_0)$ has $\sigma_T=0$, and height-0 elements are in $F_1=\{t_1\}$ with known sign. So extra terms either vanish by minimality or equal $\pm\varepsilon_1$ (known). The equations remain solvable and always yield $\sigma_T(t'\pm a_1)\neq 0$: bi-infinite chain. ✓

For $\mu<0$: extra $j\geq 5$ terms in the equations at height $h_0-c_j$ where $c_j=1+\mu+\ldots$. Since $\mu<0$, $c_5=1+\mu+\beta(a_4)+\beta(a_5-a_5)+\ldots$ The heights $c_j$ are still positive (by convexity: all polygon vertices have $h\geq 0$ and $c_j=h(x_j)-h(x_1)=h(x_j)\geq 0$; and $c_j>0$ for $j\geq 2$ since $x_j\neq x_1$ and $h(x_1)=0$ is the minimum). So $h_0-c_j<h_0$ for $c_j>0$: extra terms at height $<h_0$. By minimality of $h_0$: these equal $0$ (if in $(0,h_0)$) or $\varepsilon_1\cdot(\text{something})$ (if $=0$). Either way, equations yield chain contradiction. ✓

\medskip

For $\mu\in\mathbb{Z}\setminus\{0\}$ ($n\geq 5$): same analysis with $\mu\neq 0$. Extra terms: $h_0-c_j$ for $j\geq 5$. Since $c_j\geq c_3=1+\mu$ and we are in the case $\mu\neq 0$: for $\mu>0$, $c_j\geq 2$, so heights $<h_0-1$; for large enough $h_0$, these are below $h_0$ and vanish by minimality. For small $h_0$: direct computation shows contradiction. ✓

In all cases, the $n\geq 5$ argument reduces to the $n=4$ argument with the extra $j\geq 5$ terms vanishing or contributing known values that do not prevent the contradiction. This completes the proof for $n\geq 5$.
\end{proof}

\bigskip

\noindent\textbf{Conclusion:} For all $n\geq 3$, having $\geq 3$ edge-direction classes leads to a contradiction (bi-infinite chain in the finite set $T$, or sign sum outside $\{-1,0,+1\}$). Therefore, $T$ can have at most 2 edge-direction classes. $\blacksquare$
